% !TeX program = xelatex
% Конспект лекции 2 (19.09, С. М. Львовский) по сканам lecture-02/photo_*.jpg,
% дополненный материалом из книг: Р. Годеман, Ю. И. Манин.
% Формат B5, шрифт PT Serif; обозначения: \mathscr{F} — предпучок/пучок,
% rho — отображения ограничения.
\documentclass[11pt,b5paper]{article}

\usepackage{fontspec}
\usepackage{polyglossia}
\setdefaultlanguage{russian}
\setotherlanguage{english}
\defaultfontfeatures{Ligatures=TeX}
\setmainfont{PT Serif}

\usepackage{geometry}
\geometry{b5paper,inner=1.5cm,outer=1.3cm,top=1.5cm,bottom=1.1cm,
  headsep=6pt}

\usepackage{amsmath,amssymb,amsthm,mathtools}
\usepackage{mathrsfs}                 % \mathscr — начертание предпучков
\usepackage{tikz}
\usetikzlibrary{arrows.meta,positioning,calc,decorations.pathreplacing,shapes.geometric,patterns}
\usepackage{tikz-cd}
\usepackage{booktabs}
\usepackage{array}
\usepackage{enumitem}
\usepackage{titlesec}
\usepackage{titletoc}
\usepackage{needspace}
\usepackage{fancyhdr}
\usepackage{indentfirst}

% Токены дизайн-системы (навык brutalism: bergside/awesome-design-skills)
\definecolor{text}{HTML}{111827}
\definecolor{primary}{HTML}{DD614C}
\definecolor{secondary}{HTML}{DAA144}
\definecolor{success}{HTML}{16A34A}
\definecolor{warning}{HTML}{D97706}
\definecolor{danger}{HTML}{DC2626}
% Цвета ссылок — как у Шейкина (тёмно-красные ссылки, зелёные цитаты)
\definecolor{sheykinlink}{rgb}{0.6,0,0}
\definecolor{sheykincite}{rgb}{0,0.6,0}
\definecolor{sheykinurl}{rgb}{0,0,0.9}

\usepackage[unicode,psdextra,linktocpage=true,colorlinks=true,
  linkcolor=sheykinlink,citecolor=sheykincite,urlcolor=sheykinurl]{hyperref}
% polyglossia (russian) делает " активным символом, что ломает подписи
% в tikz-cd (\ar[r,"f"]); нам " не нужен — в тексте используются «».
\AtBeginDocument{\catcode`\"=12\relax\color{text}}

%------------------------------------------------------------------
% Обозначения
%------------------------------------------------------------------
\newcommand{\et}{\text{ét}}
\newcommand{\Hom}{\operatorname{Hom}}
\newcommand{\Sh}{\mathrm{Sh}}
\newcommand{\PSh}{\mathrm{PSh}}
\newcommand{\id}{\mathrm{id}}
\newcommand{\Shv}[1]{\mathscr{#1}}   % предпучок/пучок
\newcommand{\Top}{\mathrm{Top}}
\newcommand{\res}{\rho}              % отображение ограничения
\newcommand{\dirlim}{\varinjlim}

\theoremstyle{plain}
\newtheorem{thminner}{Утверждение}[section]
\newtheorem{leminner}[thminner]{Лемма}
\newtheorem{factinner}[thminner]{Факт}
\newtheorem{corinner}[thminner]{Следствие}
\theoremstyle{definition}
\newtheorem{defninner}[thminner]{Определение}
\newtheorem{exminner}[thminner]{Пример}
\newtheorem{assinner}[thminner]{Предположение}
\theoremstyle{remark}
\newtheorem{reminner}[thminner]{Замечание}

% Теоремные окружения набираются целиком в коробке: заголовок никогда
% не оторвётся от текста, блок не рвётся на границе страниц.
\newsavebox{\thmbox}
\newcommand{\thmbeg}{%
  \par\smallskip
  \begin{lrbox}{\thmbox}%
  \begin{minipage}{\linewidth}}
\newcommand{\thmend}{%
  \end{minipage}%
  \end{lrbox}%
  \noindent\usebox{\thmbox}%
  \par\smallskip}
\newenvironment{thm}{\thmbeg\begin{thminner}}{\end{thminner}\thmend}
\newenvironment{lem}{\thmbeg\begin{leminner}}{\end{leminner}\thmend}
\newenvironment{fact}{\thmbeg\begin{factinner}}{\end{factinner}\thmend}
\newenvironment{cor}{\thmbeg\begin{corinner}}{\end{corinner}\thmend}
\newenvironment{defn}{\thmbeg\begin{defninner}}{\end{defninner}\thmend}
\newenvironment{exm}{\thmbeg\begin{exminner}}{\end{exminner}\thmend}
\newenvironment{ass}{\thmbeg\begin{assinner}}{\end{assinner}\thmend}
\newenvironment{rem}{\thmbeg\begin{reminner}}{\end{reminner}\thmend}

% Блоки-дополнения: серый текст без рамок (стиль Шейкина)
\newenvironment{bookbox}{\par\smallskip\begingroup\color{gray}\noindent
  \ignorespaces}{\par\endgroup\smallskip}
\newcommand{\booksrc}[1]{%
  \noindent{\bfseries\small Дополнение из книги:}\par\smallskip
  \noindent#1\par\smallskip}

\newcommand{\sectionsrc}[1]{\par\smallskip\noindent
  {\color{gray}\footnotesize По сканам: #1}\par\medskip}

%------------------------------------------------------------------
% Нумерация рисунков и диаграмм (не «плавающие»: стоят на своём месте)
%------------------------------------------------------------------
\newcounter{ris}[section]
\renewcommand{\theris}{\thesection.\arabic{ris}}
\newcounter{dgm}[section]
\renewcommand{\thedgm}{\thesection.\arabic{dgm}}
\newcommand{\risCap}{}
\newcommand{\dgmCap}{}

\newsavebox{\risbox}
\newsavebox{\dgmbox}
\newenvironment{risunok}[1]{%
  \par\medskip\refstepcounter{ris}%
  \renewcommand{\risCap}{#1}%
  \par\centering\begingroup
  \begin{lrbox}{\risbox}%
  \begin{minipage}{\linewidth}%
  \centering}
 {\par\smallskip{\small Рис.~\theris.~\risCap}%
  \end{minipage}%
  \end{lrbox}%
  \usebox{\risbox}%
  \endgroup\par\medskip}

\newenvironment{diagram}[1]{%
  \par\medskip\refstepcounter{dgm}%
  \renewcommand{\dgmCap}{#1}%
  \par\centering\begingroup
  \begin{lrbox}{\dgmbox}%
  \begin{minipage}{\linewidth}%
  \centering}
 {\par\smallskip{\small Диаг.~\thedgm.~\dgmCap}%
  \end{minipage}%
  \end{lrbox}%
  \usebox{\dgmbox}%
  \endgroup\par\medskip}

% стиль для диаграмм tikz-cd
\tikzset{cd style/.style={column sep=1.7em,row sep=1.6em,
  every node/.style={font=\small}}}

% единый стиль иллюстраций (рисунков); цвета — токены brutalism
\tikzset{
  fig base/.style   = {line width=0.9pt,draw=text},
  fig open/.style   = {line width=0.7pt,dashed,draw=primary},
  fig open2/.style  = {line width=0.7pt,dashed,draw=success},
  fig fiber/.style  = {line width=0.55pt,dotted,draw=gray!75!black},
  fig sect/.style   = {line width=1.4pt,draw=danger},
  fig leafA/.style  = {line width=1.4pt,draw=primary},
  fig leafB/.style  = {line width=1.4pt,draw=success},
  fig dot/.style    = {circle,fill=text,inner sep=0pt,minimum size=4.6pt},
  fig note/.style   = {font=\footnotesize,inner sep=1.5pt},
}

\setlength{\emergencystretch}{3em}
\tolerance=3000
\setlength{\hfuzz}{2pt}
\setlist{nosep,leftmargin=*}
\setlist[enumerate,1]{label=[\arabic*],leftmargin=*}
\newlist{conditions}{enumerate}{2}
\setlist[conditions,1]{
	label=(\roman{conditionsi}),
	leftmargin=1cm,
	rightmargin=.5cm
}
\setlist[conditions,2]{
	label=(\roman{conditionsi}.\alph{conditionsii}),
	leftmargin=1cm,
	rightmargin=.5cm
}
\newlist{bullets}{itemize}{3}
\setlist[bullets,1]{label=$\circ$,leftmargin=1cm,rightmargin=.5cm}
\setlist[bullets,2]{label=--,leftmargin=1cm,rightmargin=.5cm}
\setlist[bullets,3]{label=$\diamond$,leftmargin=1cm,rightmargin=.5cm}

%--- Заголовки: \section простые (как у Шейкина), подзаголовки ---
%--- runin (жирный / курсив), как обговорено ------------------------
\titleformat{\section}[block]
  {\normalfont\Large\bfseries}
  {\thesection.}{.5em}{}
\titleformat{name=\section,numberless}[block]
  {\normalfont\Large\bfseries}
  {}{0pt}{}
\titlespacing*{\section}{0pt}{16pt}{8pt}
\titleformat{\subsection}[runin]
  {\normalfont\bfseries}
  {\thesubsection.}{.5em}{}
  [.]\titlespacing{\subsection}{0pt}{1.5ex plus .1ex minus .2ex}{1ex}
\titleformat{name=\subsubsection,numberless}[runin]
  {\normalfont\itshape}{}{0pt}{}[.]
\titlespacing{name=\subsubsection,numberless}{0pt}
  {1.5ex plus .1ex minus .1ex}{1ex}
% Заголовок не остаётся один внизу страницы: \Needspace до вызова
% (внутри формата titlesec он не срабатывает).
\let\brutsection\section
\renewcommand{\section}{\Needspace{7\baselineskip}\brutsection}
\let\brutsubsection\subsection
\renewcommand{\subsection}{\Needspace{5\baselineskip}\brutsubsection}

\pagestyle{fancy}
\fancyhf{}
\renewcommand{\headrulewidth}{0.4pt}
\renewcommand{\footrulewidth}{0pt}
\fancyhead[L]{\footnotesize Лекция~2. Прокачивание пучка}
\fancyhead[R]{\footnotesize\thepage}

\makeatletter
% Оглавление: единая колонка номеров (2.6em), заголовки с одного края,
% номера страниц — по правому краю текстовой полосы.
\titlecontents{section}
  [0pt]
  {\addvspace{10pt}\normalfont\bfseries}
  {\contentslabel[\thecontentslabel]{2.6em}}
  {}
  {\hfill\contentspage}
  [\vspace{3pt}{\color{text}\hrule height0.4pt}]
\titlecontents{subsection}
  [0pt]
  {\addvspace{2pt}\normalfont}
  {\contentslabel[\thecontentslabel]{2.6em}}
  {}
  {\titlerule*[0.45em]{.}\contentspage}
\makeatother
%------------------------------------------------------------------

\begin{document}

\thispagestyle{empty}
\vspace*{1.2cm}
\begin{center}
{\large Конспект лекции~2, 19.09\par}
\vspace{0.9cm}
{\Large\bfseries Пространство ростков $\et(\Shv{F})$, пучок
 $\Shv{F}^{\#}$\\[4pt]
 и изоморфизм
 $\Hom(\Shv{F},\Shv{G})\simeq\Hom(\Shv{F}^{\#},\Shv{G})$\par}
\vspace{0.9cm}
{\large С.~М.~Львовский --- курс теории пучков\par}
\vspace{2.2cm}
{\small Конспект по сканам lecture-02/photo\_1\,\ldots\,photo\_7
 (рукописная фолиация стр.~9---15); текст дополнен материалом из книг
 Годемана и Манина (отмечено <<Дополнение из книги>>). Формат~B5.\par}
\end{center}
\bigskip

\tableofcontents

%==================================================================
\section{План лекции, обозначения и соответствие сканам}
%==================================================================

\subsection{О чём эта лекция}
В лекции~1 предпучки и пучки были определены аксиомами (F1) и (F2),
построено пространство ростков, эскизно намечен <<прокачивание>>
(sheafification) предпучка. Во второй лекции эта конструкция доводится
до конца и получает функторное оправдание. Если коротко, лекция
отвечает на три вопроса.

\begin{conditions}
\item \textbf{Откуда берётся пучок из предпучка.} Из предпучка
 $\Shv{F}$ строится топологическое пространство $\et(\Shv{F})$ с
 локальным гомеоморфизмом $\pi\colon\et(\Shv{F})\to X$; пучком
 называется пучок непрерывных сечений $\pi$ --- он обозначается
 $\Shv{F}^{\#}$.
\item \textbf{Почему это действительно пучок.} Проверяются аксиомы
 (F1) и (F2) для $\Shv{F}^{\#}$; разбирается пример постоянного
 предпучка.
\item \textbf{Почему «прокачивание» не теряет информации.}
 Доказывается главное утверждение
 \[
   \Hom(\Shv{F},\Shv{G})\simeq\Hom(\Shv{F}^{\#},\Shv{G})
 \]
 естественно по $\Shv{F}$ и по $\Shv{G}$ при пучке $\Shv{G}$.
 Из этого следует, что $\Shv{F}\mapsto\Shv{F}^{\#}$ задано однозначно
 (с точностью до изоморфизма), и что <<прокачивание>> можно делать
 внутри готового пучка: секции $\Shv{F}^{+}$ дают тот же результат,
 $\Shv{F}^{+}\simeq\Shv{F}^{\#}$.
\end{conditions}

Последняя часть лекции --- переход к абелевой алгебре пучков: ядро,
образ (образ надо <<прокачивать>>) и точные последовательности.

\subsection{Обозначения}
В этом конспекте для предпучков и пучков используется начертание
$\Shv{F},\Shv{G},\Shv{H},\Shv{P}$ (mathscr), а отображения ограничения
пишутся через $\rho$. Остальные обозначения собраны в таблице.

\begin{center}
\small
\renewcommand{\arraystretch}{1.25}
\begin{tabular}{@{}p{0.30\textwidth}p{0.63\textwidth}@{}}
\toprule
Обозначение & Что означает \\
\midrule
$X$ & топологическое пространство (база) \\
$U,V,W,U_i$ & открытые подмножества $X$ \\
$\Top_X$ & категория открытых множеств $X$: объекты --- открытые,
 морфизмы --- включения \\
$\Shv{F},\Shv{G}$ & предпучок (а если сказано <<пучок>> --- пучок) на $X$ \\
$\Shv{F}(U)$ & сечения $\Shv{F}$ над $U$; иногда $\Gamma(U,\Shv{F})$ \\
$\res^{V}_{U}\colon\Shv{F}(V)\to\Shv{F}(U)$ &
 отображение ограничения при $U\subset V$; коротко
 $s|_{V}:=\res^{V}_{U}(s)$ \\
$\res_i:=\res^{U_i}_{U}$, $\res_x:=\res^{V_x}_{U}$ &
 те же ограничения с короткой записью \\
$\Shv{F}_x$ & слой (пространство ростков) над точкой $x$ \\
$s_x$, $\tilde s(x)$ & росток (зародыш) сечения $s$ в точке $x$ \\
$\et(\Shv{F})$ & пространство ростков (<<этажное>> пространство) \\
$\pi\colon\et(\Shv{F})\to X$ & проекция, $\pi|_{\Shv{F}_x}\equiv x$ \\
$\Shv{F}^{\#}$ & пучок сечений $\pi$, <<прокачивание>> $\Shv{F}$ \\
$i\colon\Shv{F}\to\Shv{F}^{\#}$ & канонический гомоморфизм предпучков \\
$\Shv{F}^{+}$ & секции пучка $\Shv{H}\supset\Shv{F}$, локально
 принадлежащие $\Shv{F}$ \\
$\PSh(X),\Sh(X)$ & категории предпучков и пучков на $X$ \\
$(\mathrm{F1}),(\mathrm{F2})$ & аксиомы единственности и склейки \\
\bottomrule
\end{tabular}
\end{center}

\subsection{Соответствие сканам}
\leavevmode\par
\begin{center}\small
\renewcommand{\arraystretch}{1.25}
\begin{tabular}{@{}p{0.20\textwidth}p{0.65\textwidth}@{}}
\toprule
Скан (стр.\ тетради) & Содержание \\
\midrule
photo\_1 (стр.~9) &
 заголовок <<Лекция 2 (19.09)>>, слово <<прокачивание>>;
 конструкция $\et(\Shv{F})$, $\pi$, $\Shv{F}^{\#}(U)=\{s\mid\pi\circ s=\id\}$;
 три свойства пучка (сложение, аксиомы, нейтральный элемент);
 пример с $A$-группой; <<фундаментальная проблема прокачивания>>;
 утверждение $\Hom(\Shv{F},\Shv{G})\simeq\Hom(\Shv{F}^{\#},\Shv{G})$
 и замечание о естественности (две диаграммы) \\
photo\_2 (стр.~10) &
 <<внезапная вставка»: проверка (F1) и (F2) для $\Shv{F}^{\#}$
 с рисунками-«листами»; введение топологии на $\et(\Shv{F})$ по
 критерию открытости; схема естественного отображения
 $x\mapsto U\mapsto\Shv{F}(U)\mapsto\tilde s(x)$ \\
photo\_3 (стр.~11) &
 проверка того, что $\pi^{-1}(U)$ открыто; диаграмма
 $F\to\Shv{F}^{\#}$ и инъективность $\Hom(\Shv{F}^{\#},\Shv{G})
 \hookrightarrow\Hom(\Shv{F},\Shv{G})$; обратное построение $\beta$
 и проблема единственности ($i(s_1)=i(s_2)$); выбор окрестностей
 $V_x$ и натурная диаграмма \\
photo\_4 (стр.~12) &
 большая коммутативная диаграмма; склейка $g_x,g_y$ аксиомой (F2)
 пучка $\Shv{G}$; вывод $\Hom(\Shv{F},\Shv{G})\simeq\Hom(\Shv{F}^{\#},
 \Shv{G})$; утверждение о критерии $F\simeq F'$; формулировка <<Факта>> \\
photo\_5 (стр.~13) &
 доказательство <<Факта>> диаграммами: $f\colon A_2\to A_1$,
 $g\colon A_1\to A_2$, вывод $gf=\id$, $fg=\id$ \\
photo\_6 (стр.~14) &
 подпредпучок $\Shv{F}\subset\Shv{H}$, секции $\Shv{F}^{+}$;
 план доказательства через <<Факт>>; построение
 $\eta_{\Shv{G}}(\varphi)=\bar\varphi$ \\
photo\_7 (стр.~15) &
 окончание: $\eta_{\Shv{G}}$ --- изоморфизм (диаграмма
 естественности); пример с $k$-формами; ядро, образ,
 точная последовательность \\
\bottomrule
\end{tabular}
\end{center}

Порядок файлов совпадает с рукописной фолиацией: photo\_1 ---
стр.~9 (продолжение тетради, начатой в лекции~1), photo\_2---10,
photo\_3---11, photo\_4---12, photo\_5---13,
photo\_6---14, photo\_7---15.

\subsection{Как пользоваться этим конспектом}
\leavevmode\par
\begin{bullets}
\item Все рисунки пронумерованы (рис.~$N.M$), все коммутативные
 диаграммы --- отдельной нумерой (диаграмма~$N.M$), все формулы, на
 которые есть ссылки, --- пронумерованы. Ссылка вида
 <<см.~диаграмму~\ref{dgm:glue}>> означает конкретную диаграмму
 с этим номером.
\item Серые блоки <<Дополнение из книги>> --- материал из книг
 Годемана и Манина;
 их можно читать отдельно, основной текст без них самодостаточен.
\item В доказательствах каждый <<шаг» сопровождается объяснением, зачем
 он нужен: если непонятен ход рассуждения, смотрите текст между
 формулами, а не только формулы.
\end{bullets}

%==================================================================
\section{Предпучок и пучок: повторение и уточнения}
%==================================================================

\subsection{Предпучок}
\leavevmode\par
\begin{defn}[предпучок]
Пусть $X$ --- топологическое пространство. \emph{Предпучком} множеств
(групп, колец $\ldots$) на $X$ называется семейство множеств
$\{\Shv{F}(U)\}$, поставленных открытым $U\subset X$, вместе с
отображениями ограничения
\[
  \res^{V}_{U}\colon\Shv{F}(V)\longrightarrow\Shv{F}(U)
  \qquad (U\subset V)
\]
такими, что
\[
  \res^{U}_{U}=\id,\qquad
  \res^{U}_{V}\circ\res^{V}_{W}=\res^{U}_{W}
  \quad\text{при}\quad U\subset V\subset W.
\]
Для $s\in\Shv{F}(V)$ пишут $s|_{U}:=\res^{V}_{U}(s)\in\Shv{F}(U)$.
\end{defn}

Короче: категория открытых множеств $\Top_X$ (объекты --- открытые
множества, морфизмы --- включения) --- это частный случай естественной
конструкции, и предпучок --- это контравариантный функтор
\[
  \Shv{F}\colon\Top_X^{\mathrm{op}}\longrightarrow \mathrm{Sets}
\]
(для абелевых групп, колец и т.\,п. --- с значениями в соответствующей
категории). Контравариантность --- это и есть то, что $\rho$ переводит
сечение на большем открытом множестве в сечение на меньшем.

\begin{rem}[что считать сечением]
Сечение --- это не функция в традиционном смысле, а <<данные, заданные
на $U$ и уменьшающиеся при уточнении>>. Типовой пример:
$\Shv{F}(U)$ --- множество непрерывных (или голоморфных, или
полиномиальных) функций на $U$, а $\rho$ --- ограничение функции.
Вся теория пучков ровно о том, что таких данных больше, чем
<<глобальных функций>>, и разрыв между ними измеряется
несвязностью $U$.
\end{rem}

\subsection{Пучок: аксиомы (F1) и (F2)}
\leavevmode\par
\begin{defn}[пучок]
Предпучок $\Shv{F}$ называется \emph{пучком}, если для любого открытого
$U$ и любого покрытия $U=\bigcup_{i\in I}U_i$ выполняются:
\begin{bullets}
\item[\textup{(F1)}] \emph{(единственность, согласованность)} если
 $s,t\in\Shv{F}(U)$ такие, что $s|_{U_i}=t|_{U_i}$ для всех $i$, то
 $s=t$;
\item[\textup{(F2)}] \emph{(склейка)} если даны $s_i\in\Shv{F}(U_i)$,
 согласованные на пересечениях,
 \[
   s_i|_{U_i\cap U_j}=s_j|_{U_i\cap U_j}\qquad\text{для всех }i,j,
 \]
 то существует $s\in\Shv{F}(U)$ с $s|_{U_i}=s_i$; причём это $s$
 определено единственным образом.
\end{bullets}
\end{defn}

Интуиция:
\begin{bullets}
\item (F1) говорит: \emph{локальное равенство --- это глобальное
 равенство}. Если два глобальных объекта неразличимы на каждом куске
 покрытия, они совпадают. То есть <<достаточно знать данные на мелких
 кусках>>.
\item (F2) говорит: \emph{согласованные локальные данные склеиваются в
 глобальные}. Обратное к (F1): <<из кусков можно собрать целое>>.
\end{bullets}
Ни одно из этих свойств не выполняется для произвольного предпучка;
именно поэтому <<пучок>> --- отдельное (и гораздо более полезное)
понятие.

\begin{bookbox}
\booksrc{Годеман, гл.~II, \S\,1, п.~1.1, с.~129 --- формулировка аксиом
 (F1), (F2) в терминах <<семейства открытых подмножеств>>; замечание:
 из (F2) в случае пустого покрытия следует, что $\Shv{F}(\varnothing)$
 --- одноэлементное множество, поэтому в (F2) достаточно рассматривать
 пары индексов с непустым пересечением.
\par\smallskip
 Манин, \S\,2.1.1---2.1.3, с.~108---110: то же определение в виде системы
 $\{\Shv{F}(U),\res^{V}_{U}\}$ и формулировка <<пучок --- пучок, если
 согласованные сечения на покрытии склеиваются и притом единственным
 образом>>. Обозначение $\Gamma(U,\Shv{F})$ для сечений.
}
\end{bookbox}

\subsection{Аксиомы пучка как точная последовательность}
Для предпучков абелевых групп аксиомы (F1), (F2) можно записать одной
строкой --- как точность последовательности. Это пригодится в конце
лекции, когда мы будем говорить о точных последовательностях пучков.

Пусть $U=\bigcup_{i\in I}U_i$ --- фиксированное покрытие. Рассмотрим
диаграмму~\ref{dgm:exact-pre}.

\begin{diagram}{Аксиомы пучка в виде точности для покрытия
 $U=\bigcup_i U_i$ (см.~также разбор ниже).\label{dgm:exact-pre}}
\begin{tikzcd}[cd style,column sep=large]
0 \ar[r] & \Shv{F}(U) \ar[r,"f"] &
\displaystyle\prod_{i\in I}\Shv{F}(U_i)
\ar[r,shift left=5pt,"\gamma"] \ar[r,shift right=5pt,"\delta"'] &
\displaystyle\prod_{i,j\in I}\Shv{F}(U_i\cap U_j)
\end{tikzcd}
\end{diagram}

Пояснение к диаграмме~\ref{dgm:exact-pre}. Читается слева
направо: стрелка $f$ <<вклеивает>> глобальное сечение в семейство
локальных, а две стрелки $\gamma$ и $\delta$ с одного множества
на другое --- это ограничение на пересечении $U_i\cap U_j$
у $i$-й и у $j$-й компоненты семейства соответственно. Точность
(равенство $\ker(\gamma-\delta)=\operatorname{Im}f$) и есть
аксиомы пучка одной строкой: «согласованные семейства --- ровно
те, что являются ограничениями глобальных сечений», --- из
неё следуют и единственность (F1), и склейка (F2); нулевой член
слева гарантирует, что $f$ ничего не теряет. Формулы компонент
--- в замечании ниже.

\begin{rem}[что означают $f,\gamma,\delta$]\label{rem:gd}
Все отображения поэлементные.
\begin{conditions}
\item Первая стрелка $f$ <<вклеивает>> глобальное сечение в семейство
 локальных:
 \begin{equation}\label{eq:f}
  f(s)=\bigl(\,\res^{U_i}_{U}(s)\,\bigr)_{i\in I}
  =\bigl(\,s|_{U_i}\,\bigr)_{i\in I}
  \in\prod_{i\in I}\Shv{F}(U_i).
 \end{equation}
\item Дальше стоят \emph{две} стрелки с одним и тем же множеством
 значений: они берут ограничение на пересечение $U_i\cap U_j$
 соответственно у $i$-й и у $j$-й компоненты семейства
 $(s_i)_{i\in I}$:
 \begin{equation}\label{eq:gd}
  \gamma\bigl((s_i)_i\bigr)_{ij}=\res^{U_i\cap U_j}_{U_i}(s_i),
  \qquad
  \delta\bigl((s_i)_i\bigr)_{ij}=\res^{U_i\cap U_j}_{U_j}(s_j).
 \end{equation}
\item Семейство $(s_i)$ называется \emph{согласованным}, если
 $\gamma\bigl((s_i)\bigr)=\delta\bigl((s_i)\bigr)$, т.\,е.\ если
 $s_i|_{U_i\cap U_j}=s_j|_{U_i\cap U_j}$ для всех $i,j$. Иными словами,
 $\gamma$ и $\delta$ --- одно и то же действие <<взять ограничение на
 пересечение>>, отнесённое к разным компонентам; их равенство и есть
 условие согласованности из (F2).
\end{conditions}
Вместо пары стрелок $\gamma,\delta$ часто пишут одну $\gamma-\delta$:
\begin{equation}\label{eq:gmd}
 (\gamma-\delta)\bigl((s_i)\bigr)_{ij}
 =\res^{U_i\cap U_j}_{U_i}(s_i)-\res^{U_j\cap U_i}_{U_j}(s_j).
\end{equation}
\end{rem}

Теперь аксиомы:
\begin{bullets}
\item \textbf{Точность в первом месте} ($f$ --- инъекция) --- это (F1):
 из $\res^{U_i}_{U}(s)=0$ для всех $i$ следует $s=0$, а для групп ---
 из равенства нулю; обобщённо: если $s|_{U_i}=t|_{U_i}$ для всех $i$,
 то $s=t$.
\item \textbf{Точность во втором месте} ($\ker(\gamma-\delta)=
 \operatorname{Im}f$) --- это (F2): слева направо --- любое согласованное
 семейство $(s_i)$ является семейством вида~\eqref{eq:f}, т.\,е.
 склеивается в глобальное $s$ (и притом единственное); справа налево ---
 глобальное сечение даёт согласованное семейство.
\item \textbf{Точность в последнем месте} не требуется: не всякое
 семейство из $\prod_{i,j}\Shv{F}(U_i\cap U_j)$ рождается из
 семейства $(s_i)$.
\end{bullets}

\begin{bookbox}
\booksrc{Манин, \S\,2.1.4, с.~110 (эта диаграмма) и \S\,2.1.5, с.~111
 (фундаментальный пример предпучка, который не пучок --- рис.~2.2 книги,
 он же рисунок~\ref{fig:circle} ниже).}
\end{bookbox}

\subsection{Пустое открытое множество}
Из (F2) для пустого покрытия $U=\varnothing$ следует, что
$\Shv{F}(\varnothing)$ --- одноэлементное множество (в абелевой
категории: нулевой элемент). Поэтому при проверке (F2) можно
рассматривать только пары индексов с непустым пересечением: если
$U_i\cap U_j=\varnothing$, условие согласованности на нём ---
тривиально, ведь ограничение в пустое множество --- единственное
отображение в однозначное множество.

\subsection{Зачем нужны именно пучки: два примера}
\leavevmode\par
\begin{exm}[склеивание двух прямых]
Пусть $X_1\simeq X_2\simeq\mathbb{R}$, $U_1\simeq U_2\simeq
\mathbb{R}\setminus\{0\}$ (две копии прямой без нуля). Склеиваем
$X_1\sqcup X_2$ по $U_1\sqcup U_2$ с помощью отображения $f$. При
$f=\id$ получается прямая с раздвоенной точкой, при $f(x)=1/x$ ---
проективная прямая (рисунок~\ref{fig:skleivanie}). Топологии
$X_1\supset U_1$ и $X_2\supset U_2$ одинаковы в обоих случаях,
и результат определяет только функция $f$. Поэтому склеивать
нужно вместе с функциями, т.\,е. работать с пучками: голая топология
разницы между $f=\id$ и $f(x)=1/x$ не видит.
\end{exm}

\begin{risunok}{Склейка двух прямых: иллюстрация к \S\,2.1 книги Манина
 (рис.~2.1).\label{fig:skleivanie}}
\resizebox{\linewidth}{!}{\begin{tikzpicture}[>=Stealth,font=\small]
  % а) Манин, рис. 2.1а: две копии прямой, стрелки --- сшиваемые пары f = id;
  % у нулей стрелки нет: 0 из первой и 0 из второй остаются двумя точками
  \draw[fig base] (-2.1,2.5) -- (2.1,2.5);
  \draw[fig base] (-2.1,1.5) -- (2.1,1.5);
  \fill (0,2.5) circle (2.4pt);
  \fill (0,1.5) circle (2.4pt);
  \node[above] at (0,2.5) {0};
  \node[below] at (0,1.5) {0};
  \foreach \x in {-1.5,-0.9,0.9,1.5}
    \draw[<->,line width=0.7pt,draw=text] (\x,1.62) -- (\x,2.38);
  \node[below,font=\footnotesize] at (0,0.75)
    {(а) $f=\id$: прямая с раздвоенной точкой};
  % б) проективная прямая
  \begin{scope}[xshift=6.6cm]
    \draw[fig base] (0,2.0) circle (1.0);
    \fill (1.0,2.0) circle (2.4pt);
    \node[below right,font=\footnotesize] at (1.05,2.1) {$\infty$};
    \node[below,font=\footnotesize] at (0,0.15)
      {(б) $f(x)=1/x$: проективная прямая};
  \end{scope}
\end{tikzpicture}}
\end{risunok}

Пояснение к рисунку~\ref{fig:skleivanie}. Слева --- результат
склейки при $f=\id$: нулевые точки двух копий не сшиваются
(стрелками показаны сшиваемые пары), получается прямая
с раздвоенной точкой: $0_1\neq 0_2$, и разнести их нельзя ---
любые окрестности пересекаются. Справа --- та же операция при
$f(x)=1/x$: <<концы>> прямых сшиваются через $1/x$, результат ---
окружность (проективная прямая), на которой обе бывшие нулевые
точки уже разделимы. Топологически $U_1,U_2$ одинаковы в обоих
случаях, различается только $f$: поэтому одной топологии
недостаточно --- нужно склеивать вместе с функциями, то есть
работать с пучками.

\begin{exm}[предпучок, не являющийся пучком]\label{exm:circle}
$X$ --- окружность, $\mathcal{O}(U)$ --- группа $\mathbb{R}$-значных
непрерывных функций на $U$, $\underline{\mathbb{Z}}$ --- <<постоянный>>
пучок. Рассмотрим предпучок
$\Shv{P}(U)=\mathcal{O}(U)/\underline{\mathbb{Z}}(U)$.
Пусть $X=U_1\cup U_2$, причём $U_1\cap U_2=V_1\sqcup V_2$ --- два
разрывных куска. Выберем $f_1\in\mathcal{O}(U_1)$, $f_2\in
\mathcal{O}(U_2)$ так, что
\[
 \res^{V_1}_{U_1}(f_1)-\res^{V_1}_{U_2}(f_2)=0,
 \qquad
 \res^{V_2}_{U_1}(f_1)-\res^{V_2}_{U_2}(f_2)=1.
\]
Классы в $\Shv{P}(V_1)$ и $\Shv{P}(V_2)$ согласованы, но глобального
сечения нет: нельзя устранить рассогласование на $V_2$, потому что
<<целое число>> нельзя поправить непрерывной функцией. Причина ---
неодносвязность окружности: аксиома (F2) здесь нарушается.
\end{exm}

\begin{risunok}{Окружность: классы $f_1$ и $f_2$ согласованы в
 $\Shv{P}$, но глобального сечения нет (Манин, рис.~2.2).
  \label{fig:circle}}
\begin{tikzpicture}[>=Stealth,font=\small]
  \draw[fig base] (0,0) circle (1.9);
  % дуги: U_1 сверху (синяя), U_2 снизу (красная), вне окружности
  \draw[fig leafA] (1.8619,1.075)
    arc[start angle=30,end angle=150,radius=2.15];
  \draw[fig sect] (1.8619,-1.075)
    arc[start angle=-30,end angle=-150,radius=2.15];
  \node[primary,font=\small] at (0,2.6) {$U_1$};
  \node[danger,font=\small] at (0,-2.6) {$U_2$};
  \node[font=\small] at (-1.25,0) {$V_1$};
  \node[font=\small] at (1.25,0) {$V_2$};
  \node[fig note,primary] at (-1.45,2.2) {$f_1$};
  \node[fig note,danger] at (1.45,-2.2) {$f_2$};
  % скачок +1 на V_2
  \draw[->,danger,line width=1pt]
    (2.2517,1.3) arc[start angle=30,end angle=-30,radius=2.6];
  \node[fig note,danger] at (3.05,0.1) {$+1$};
\end{tikzpicture}
\end{risunok}

Пояснение к рисунку~\ref{fig:circle}. Окружность распадается на
два открытых куска: $U_1$ (верхняя дуга, синяя) и $U_2$ (нижняя,
красная); их пересечение --- два разрывных куска $V_1$ слева и
$V_2$ справа. Синяя и красная линии --- это <<значения>> функций
$f_1$ и $f_2$ на соответствующих кусках. На $V_1$ концы дуг
совпадают (разность равна $0$), а на $V_2$ расходятся на $1$
(стрелка $+1$): согласованность есть, но глобального сечения
нет, потому что <<поправить>> скачок непрерывной функцией
нельзя --- в этом и состоит нарушение (F2).

%==================================================================
\section[Ростки и слои: $\Shv{F}_x$]{Ростки и слои: $\Shv{F}_x$}
%==================================================================

\subsection{Определение}
\leavevmode\par
\begin{defn}[слой, росток]
Пусть $\Shv{F}$ --- предпучок на $X$, $x\in X$. \emph{Слоем}
(пространством ростков над $x$) называется прямой предел
\begin{equation}\label{eq:stalk}
 \Shv{F}_x=\dirlim_{U\ni x}\Shv{F}(U),
\end{equation}
где $U$ пробегает открытые окрестности точки $x$, упорядоченные
включением (по контравариантности $\Shv{F}$ предел --- прямой над
категорией окрестностей).
\end{defn}

Элементы $\Shv{F}_x$ называются \emph{ростками} (зародышами). У
предельной конструкции есть наглядное описание, которое удобно
использовать на практике: элементы $\Shv{F}_x$ --- это классы
эквивалентности пар $(U,s)$, где $U\ni x$ открыто, $s\in\Shv{F}(U)$, а
эквивалентность
\[
 (U,s)\sim(V,t)
 \quad\Longleftrightarrow\quad
 \exists\, W\ni x,\ W\subset U\cap V:\ \res^{W}_{U}(s)=\res^{W}_{V}(t).
\]
Класс пары $(U,s)$ обозначают $s_x$ или $\tilde s(x)$: <<росток
сечения $s$ в точке $x$>>. Так записывается вычислительное правило:

\begin{quote}
\itshape Чтобы найти $s_x$, возьмите любую (достаточно малую)
окрестность $U$ точки $x$ и посмотрите, чем является $s$ на ней;
если на более малой окрестности $s$ совпал с другим сечением ---
ростки совпадают.
\end{quote}

\begin{exm}
Пусть $\Shv{F}(U)$ --- непрерывные $\mathbb{R}$-значные функции на
$U$, $\rho$ --- ограничение. Тогда $\Shv{F}_x$ --- это <<локальные
функции в точке $x$>>: два представителя $(U,f)$ и $(V,g)$ дают один
и тот же росток, если $f=g$ на некоторой окрестности $x$. Это
классическая конструкция зародыша функции из математического анализа.
\end{exm}

\begin{exm}[постоянный предпучок]
Если $\Shv{F}(U)=A$ для всех непустых $U$ и все $\rho$ ---
тождественные, то в пределе ничего не <<схлопывается>>:
$\Shv{F}_x=A$ для всякой $x$. Подробнее --- в
\S\,\ref{sec:const}.\end{exm}

\begin{risunok}{Как строится росток: вложенные окрестности
 $U\supset V\supset W\ni x$ и цепочка ограничений;
 росток --- класс эквивалентности <<стопки>> сечений.}
\begin{tikzpicture}[>=Stealth,font=\small]
  \draw[fig base] (0,0) ellipse (2.5 and 1.0);
  \draw[fig open] (0,0) ellipse (1.6 and 0.62);
  \draw[danger,line width=1pt,dotted] (0,0) ellipse (0.8 and 0.3);
  \fill (0,0) circle (2.8pt);
  \node[below right] at (0.12,-0.17) {$x$};
  \node[left] at (-2.5,0.42) {$U$};
  \node[left] at (-1.6,-0.45) {$V$};
  \node[above right] at (0.7,0.42) {$W$};
  % цепочка ограничений
  \node at (4.4,1.2)  {$\Shv{F}(U)$};
  \node at (4.4,0)    {$\Shv{F}(V)$};
  \node at (4.4,-1.2) {$\Shv{F}(W)$};
  \node at (6.7,-1.2) {$\Shv{F}_x$};
  \draw[->] (4.4,0.9) -- (4.4,0.3)
    node[midway,right,font=\footnotesize] {$\res$};
  \draw[->] (4.4,-0.9) -- (4.4,-0.3)
    node[midway,right,font=\footnotesize] {$\res$};
  \draw[->] (5.35,-1.2) -- (6.05,-1.2);
\end{tikzpicture}
\end{risunok}

Пояснение к рисунку. Слева --- вложенные окрестности точки $x$:
чёрная сплошная $U$, синяя пунктирная $V$, красная точками $W$;
в центре --- сама точка $x$. Справа --- что делает с этим
предпучок: сечение на $U$ ограничивается до $V$, затем до $W$
(стрелки $\rho$), и в пределе по такой цепочке получается росток
$\Shv{F}_x$. Росток и <<стопка>> сечений --- одно и то же:
два набора дают один и тот же росток, если они совпадают на
какой-то более мелкой окрестности.

\subsection{Почему слои важны}
Из лекции~1 мы помним главное правило: \emph{пучок определяется своими
слоями}, а точнее --- гомоморфизм пучков, который поэлементно (в каждом
слое) является изоморфизмом, есть изоморфизм. Это следствие того, что
аксиома (F1) позволяет <<восстановить>> глобальное сечение по его
росткам. Поэтому <<прокачивание>>, которое мы построим, ничего не
потеряет: слои $\Shv{F}_x$ и $\Shv{F}^{\#}_x$ будут совпадать.

\begin{bookbox}
\booksrc{Манин, \S\,2.1.6, с.~112 --- определение слоя
 $\Shv{P}_x=\dirlim \Shv{P}(U)$ и ростка.
\par\smallskip
 Годеман, гл.~II, \S\,1, п.~1.2, с.~130---133 --- накрывающее
 пространство, связанное с пучком (в оригинале \emph{espace étalé};
 переводчик предпочёл <<накрывающее пространство>>).
}
\end{bookbox}

%==================================================================
\section[Пространство ét(F) и его топология]
 {Пространство $\et(\Shv{F})$ и его топология}
%==================================================================

\subsection{Множество ростков}
\leavevmode\par
\begin{defn}
Положим
\begin{equation}\label{eq:etale-set}
 \et(\Shv{F}):=\coprod_{x\in X}\Shv{F}_x
\end{equation}
--- разность множеств слоёв, и обозначим через
$\pi\colon\et(\Shv{F})\to X$ проекцию, отображающую слой $\Shv{F}_x$
в точку $x$. Тогда $\pi^{-1}(x)=\Shv{F}_x$.
\end{defn}

Пока это просто множество: точки пространства --- это пары
<<точка базы + росток в ней>>. Осталось дать топологию.

\subsection{Топология}
\leavevmode\par
\begin{defn}[топология на $\et(\Shv{F})$]
На $\et(\Shv{F})$ вводится наиболее тонкая топология, при которой все
отображения
\[
 s\colon U\longrightarrow\et(\Shv{F}),\qquad
 s(x)=s_x\ \ (s\in\Shv{F}(U)),
\]
непрерывны.
\end{defn}

Здесь стоит понять, что именно определено. У нас есть семейство
отображений $s\colon U\to\et(\Shv{F})$ (одно для каждого открытого $U$
и каждого $s\in\Shv{F}(U)$). Требуем, чтобы они все были непрерывными,
и берём наиболее тонкую топологию, которая это обеспечивает: множество
$G\subset\et(\Shv{F})$ объявляется открытым тогда и только тогда, когда
прообраз $s^{-1}(G)$ открыт для всякой такой секции $s$. Семейство таких
множеств образует топологию. Эквивалентная, <<вычислимая>> формулировка
--- критерий открытости.

\begin{thm}[критерий открытости]\label{thm:open}
Подмножество $G\subset\et(\Shv{F})$ открыто тогда и только тогда, когда
для всякого открытого $U\subset X$ и всякого $s\in\Shv{F}(U)$ множество
\begin{equation}\label{eq:crit}
 \{\,x\in U:\ s_x\in G\,\}=s^{-1}(G)\subset U
\end{equation}
открыто в $U$.
\end{thm}

\begin{proof}
($\Leftarrow$) Обозначим через $\tau$ семейство подмножеств $G$,
удовлетворяющих условию. Проверим, что $\tau$ --- топология.
$\varnothing$ и $\et(\Shv{F})$ проходят условие очевидно. Объединение:
если $G=\bigcup_\alpha G_\alpha$, то
$s^{-1}(G)=\bigcup_\alpha s^{-1}(G_\alpha)$ --- объединение открытых.
Пересечение: $s^{-1}(G_1\cap G_2)=s^{-1}(G_1)\cap s^{-1}(G_2)$ ---
пересечение открытых. Значит $\tau$ --- топология, и по построению
все $s$ в ней непрерывны. Любая другая топология, делающая все $s$
непрерывными, состоит только из множеств, удовлетворяющих условию
\eqref{eq:crit}, то есть содержится в $\tau$. Поэтому $\tau$ --- наиболее
тонкая такая топология.

($\Rightarrow$) Пусть топология на $\et(\Shv{F})$ выбрана так, как в
определении, и $G$ открыто. Тогда $s^{-1}(G)$ --- прообраз открытого
множества при непрерывном отображении, значит открыт.
\end{proof}

\sectionsrc{формулировка критерия --- photo\_2 (стр.~10) и
 photo\_3 (стр.~11): <<подмножество $G\subset\et(\Shv{F})$
 открыто $\Leftrightarrow$ множество всех ростков, в которые попадают
 локации сечения, открыто»; на стр.~11 отдельно проверяется, что
 прообраз $\pi^{-1}(U)$ открыт: <<$H=\pi^{-1}(U)$ состоит из
 всевозможных ростков $s_x$, $x\in U$, $s\in\Shv{F}(U)$, тогда оно
 открыто, т.к. $U$ открыто».}

\subsection{\texorpdfstring{$\pi$}{} --- локальный гомеоморфизм}
\leavevmode\par
\begin{thm}\label{thm:localhomeo}
Отображение $\pi\colon\et(\Shv{F})\to X$ является локальным
гомеоморфизмом: для каждой точки $\xi\in\et(\Shv{F})$ найдётся открытая
окрестность, на которой $\pi$ --- гомеоморфизм на открытое подмножество
$X$.
\end{thm}

\begin{proof}
Сначала проверим непрерывность $\pi$. Для любого открытого $O\subset X$
и всякой секции $t\in\Shv{F}(V)$ имеем
$t^{-1}(\pi^{-1}(O))=V\cap O$, открытое в $V$; по критерию
открытости $\pi^{-1}(O)$ открыто.
Пусть $\xi=s_x\in\Shv{F}_x$ для некоторого $s\in\Shv{F}(U)$.
Рассмотрим отображение $\widetilde s:U\to\et(\Shv{F})$,
$y\mapsto s_y$. Оно непрерывно по определению топологии и
$\pi\circ\widetilde s=\id_U$, поэтому инъективно. Покажем, что его
образ открыт. Для любой секции $t\in\Shv{F}(V)$ по критерию
\ref{thm:open} имеем
\[
 t^{-1}(\widetilde s(U))=\{y\in V\cap U:t_y=s_y\}.
\]
Если $t_y=s_y$, то по определению равенства ростков существует открытая
окрестность $W_y\subset V\cap U$ точки $y$, на которой
$t|_{W_y}=s|_{W_y}$. Значит, указанное множество открыто в $V$;
следовательно, $\widetilde s(U)$ открыто в $\et(\Shv{F})$.
Ограничение $\pi|_{\widetilde s(U)}$ непрерывно, а его обратное ---
$\widetilde s$, поэтому это гомеоморфизм на открытое $U$. Так как
$\xi\in\widetilde s(U)$, $\pi$ является локальным гомеоморфизмом.
\end{proof}

Следствие: $(\et(\Shv{F}),\pi)$ --- étale-пространство над $X$:
у каждой его точки есть окрестность, отображающаяся гомеоморфно на
открытое подмножество базы. В общем случае $\pi$ не обязано быть
накрытием и $\pi^{-1}(U)$ не обязано быть произведением
$U\times(\text{дискретное множество})$.

\begin{bookbox}
\booksrc{Годеман, гл.~II, \S\,1, п.~1.1, с.~129---130 --- язык
 расслоений. \emph{Расслоённым пространством} с базой $X$ называют
 пару $(E,p)$, где $E$ --- топологическое пространство, а
 $p\colon E\to X$ непрерывно; множество $E(x)=p^{-1}(x)$ обычно
 называют слоем над точкой $x$ (примечание ред., с.~129).
 \emph{Сечением} над подмножеством $M\subset X$ называется
 непрерывное $s\colon M\to E$ с $p\circ s=\id$; множество $E(U)$
 всех сечений над открытым $U$, с ограничениями, --- это пучок
 (пучок ростков сечений расслоённого пространства). Примечание
 переводчика там же: автор писал \emph{espace étalé}, русский
 перевод предпочёл <<накрывающее пространство>> --- в более
 широком, чем обычно, смысле; \emph{ét} в $\et(\Shv{F})$ ---
 от этого же термина.}
\end{bookbox}

\begin{center}
\normalsize
\setlength{\tabcolsep}{6pt}
\renewcommand{\arraystretch}{1.25}
\begin{tabular}{@{}>{\raggedright\arraybackslash}p{0.30\textwidth}
 >{\raggedright\arraybackslash}p{0.46\textwidth}
 >{\raggedright\arraybackslash}p{0.15\textwidth}@{}}
\toprule
\textbf{Из Годемана (расслоения)} &
\textbf{В этой лекции (пучки)} &
\textbf{Где} \\
\midrule
пара $(E,p)$, база $X$ & $\pi\colon\et(\Shv{F})\to X$ & \S\,4.1 \\
слой $E(x)=p^{-1}(x)$ & слой ростков $\Shv{F}_x$ & \S\,3.1 \\
сечение $s$, $p\circ s=\id$ & элемент $\Shv{F}^{\#}(U)$ &
 \S\,4.4--4.5 \\
$p$ --- локальный гомеоморфизм (накрывающее пространство) &
 $\pi$ --- локальный гомеоморфизм & \S\,4.3 \\
пучок ростков сечений $U\mapsto E(U)$ & $\Shv{F}^{\#}$ & \S\,5 \\
изоморфизм расслоений & изоморфизм пучков & \S\,8 \\
\bottomrule
\end{tabular}
\end{center}
\par\smallskip

\subsection{Сечение как непрерывный разрез}
Рисунок~\ref{fig:etale} показывает конструкцию целиком: над каждой
точкой $X$ висит свой слой, а сечение --- это <<разрез>>, который над
каждой точкой выбирает росток и делает это непрерывно.

\begin{risunok}{Проекция $\pi\colon\et(\Shv{F})\to X$: над каждой
 точкой --- слой ростков; сечение $s\in\Shv{F}^{\#}(U)$ --- непрерывный
 разрез, проецирующийся в себя ($\pi\circ s=\id$).\label{fig:etale}}
\begin{tikzpicture}[>=Stealth,font=\small]
  % база X с открытым U
  \draw[fig base] (0,0) ellipse (3.6 and 0.9);
  \node[below right] at (3.3,-0.55) {$X$};
  \begin{scope}
    \clip (0,0) ellipse (3.6 and 0.9);
    \fill[primary!10] (-1.8,0) ellipse (1.9 and 0.85);
  \end{scope}
  \draw[fig open] (-1.8,0) ellipse (1.9 and 0.85);
  \node[primary,below] at (-1.8,-1.25) {$U$};
  % этажёрка ét(F) — без рамки
  \node[anchor=north west] at (-4.15,5.15) {$\et(\Shv{F})$};
  % волокна и ростки
  \draw[fig fiber] (-2.6,0.63) -- (-2.6,4.9);
  \draw[fig fiber] (-0.9,0.88) -- (-0.9,4.9);
  \draw[fig fiber] ( 1.4,0.83) -- ( 1.4,4.9);
  \draw[fig fiber] ( 3.0,0.50) -- ( 3.0,4.9);
  \foreach \x/\y in {-2.6/1.7,-2.6/3.8,-0.9/2.2,-0.9/4.2,
                     1.4/1.9,1.4/3.9,3.0/2.3,3.0/4.3}
    \fill (\x,\y) circle (2.4pt);
  \foreach \x/\y in {-2.6/2.6,-0.9/3.2}
    \fill[danger] (\x,\y) circle (2.6pt);
  % сечение над U
  \draw[fig sect,smooth] plot coordinates
    {(-3.45,2.35) (-2.6,2.6) (-0.9,3.2) (0.15,3.05)};
  \node[danger] at (-0.3,3.62) {$s$};
  % слой над точкой x
  \draw[decorate,decoration={brace,mirror,raise=4pt}]
    (-2.6,1.7) -- (-2.6,3.8);
  \node[rotate=90,anchor=south,font=\footnotesize] at (-3.4,2.75)
    {$\pi^{-1}(x)=\Shv{F}_x$};
  % проекция
  \draw[->,thick] (4.75,5.0) -- (4.75,0.35);
  \node at (4.98,2.7) {$\pi$};
\end{tikzpicture}
\end{risunok}

Пояснение к рисунку~\ref{fig:etale}. Прямоугольник сверху ---
пространство ростков $\et(\Shv{F})$, овал снизу --- база $X$.
Вертикальные пунктирные линии --- волокна $\pi^{-1}(x)$: над
каждой точкой $x$ висит свой слой ростков (столбик чёрных точек),
скобкой слева показано волокно над одной точкой,
$\pi^{-1}(x)=\Shv{F}_x$. Синим выделено открытое $U\subset X$;
красная кривая $s$ --- сечение над $U$: для каждой точки $x\in U$
она выбирает росток (красные точки на волокнах), причём выбор
меняется непрерывно, поэтому кривая не рвётся. Стрелка $\pi$
--- проекция вниз, и $\pi\circ s=\id$ означает: каждый выбранный
росток лежит ровно над своей точкой.

\subsection{Сечение и мульти-секция}
Непрерывность $s$ означает, что выбор ростков по точкам $x\in U$
происходит <<плавно>>: множество $\{s(x):x\in U\}$ открыто в
$\et(\Shv{F})$. Полезно противопоставить этому произвольный выбор
ростков (рисунок~\ref{fig:multisection}): можно для каждой точки $U$
произвольно вытащить из слоя какой угодно росток --- получится
<<мульти-секция>>, но, вообще говоря, не сечение.

\begin{risunok}{Сечение --- непрерывный выбор ростков (красное); синий
 произвольный выбор ростков по точкам $U$ сечением не является.
 Формально $\Shv{F}^{\#}(U)=\{s\mid(\pi\circ s)(U)=\id_U\}$.
  \label{fig:multisection}}
\begin{tikzpicture}[>=Stealth,font=\small]
  \draw[fig base] (0,0) ellipse (3.4 and 0.75);
  \node[below] at (0,-0.8) {$U$};
  \foreach \x/\top in {-2.6/0.483,-1.3/0.693,0/0.75,1.3/0.693,2.6/0.483}
    \draw[fig fiber] (\x,\top) -- (\x,3.5);
  \foreach \x in {-2.6,-1.3,0,1.3,2.6}
    \foreach \y in {1.5,2.4,3.2} \fill (\x,\y) circle (2.4pt);
  % непрерывное сечение
  \draw[fig sect,smooth] plot coordinates
    {(-3.0,1.2) (-2.6,1.5) (-1.3,2.4) (0,3.2) (1.3,2.4) (2.6,1.5) (3.0,1.2)};
  \node[danger,above] at (0,3.55) {$s\in\Shv{F}^{\#}(U)$};
  % произвольная мульти-секция
  \fill[warning] (-2.6,2.4) circle (3.2pt);
  \fill[warning] (0,1.5) circle (3.2pt);
  \fill[warning] (2.6,3.2) circle (3.2pt);
  \draw[warning,dashed] (-2.6,2.4) -- (0,1.5) -- (2.6,3.2);
  \node[warning,font=\footnotesize] at (0,-1.95)
    {произвольный выбор ростков: мульти-секция, не сечение};
\end{tikzpicture}
\end{risunok}

Пояснение к рисунку~\ref{fig:multisection}. Внизу --- открытое
$U\subset X$, вверх от него идут волокна: столбики ростков над
точками $U$. Красная кривая --- сечение: она соединяет ростки в
непрерывный <<разрез>> и проецируется в $U$ «на себя». Синие
точки --- произвольный выбор ростка над каждой точкой $U$:
соединённые пунктиром, они скачут между уровнями, выбор
не непрерывен, поэтому сечением не является.

\begin{bookbox}
\booksrc{Годеман, гл.~II, \S\,1, п.~1.3, с.~133 --- сечения
 определяются над произвольным подмножеством $M\subset X$, не
 только над открытым. При этом (F1) остаётся осмысленной и верной
 без открытости, а (F2) --- нет: иначе любое отображение
 $s\colon X\to E$ с $p\circ s=\id$ было бы непрерывным, ведь
 $X$ распадается на одноточечные множества. Теорема~1.3.1 там же:
 склейка сечений работает на локально конечном \emph{замкнутом}
 покрытии.\par\smallskip
 Манин, \S\,1.10, с.~63 --- расслоенное произведение;
 \S\,1.12, с.~77---79 (определение~1.12.9, с.~79): векторное
 расслоение --- это семейство векторных пространств, локально
 тривиальное в каждой точке. Это прямой прототип будущего пучка
 модулей: <<расслоение>> в алгебраизированном виде.}
\end{bookbox}

\subsection{Нехаусдорфовость}
\leavevmode\par
\begin{rem}[нехаусдорфовость]
Даже если $X$ хаусдорфово, $\et(\Shv{F})$, вообще говоря, не
хаусдорфово. По Годеману (Замечание 1.2.4, с.~133): если
$\et(\Shv{F})$ отделимо, то два сечения $s,t$ над открытым $U$,
совпадающие в некоторой точке, равны во всей связной компоненте,
т.\,е. в пучке действует <<принцип аналитического продолжения>>.
Например, пучок ростков голоморфных функций на комплексном
многообразии отделим, а пучок ростков непрерывных функций --- нет.
\end{rem}

\begin{risunok}{Пример нехаусдорфова $\et(\Shv{F})$: $f_1$ и $f_2$
 совпадают на $[0,\tfrac12]$ и расходятся дальше. Ростки над точкой
 $\tfrac12$ различны, но любые их окрестности пересекаются, так как
 листы соединяются слева (иллюстрация к Манину, рис.~2.3, и к заметке
 Годемана 1.2.4).\label{fig:nonhausdorff}}
\begin{tikzpicture}[>=Stealth,font=\small]
  % левый график: f_1 и f_2
  \draw[->] (-0.3,0) -- (4.45,0);
  \draw[->] (0,-0.3) -- (0,3.15);
  \draw[fig sect,dashed] (0.3,0.15) -- (4.15,0.15);
  \draw[fig leafA,smooth] plot coordinates
    {(0.3,0.15) (1.7,0.15) (2.5,0.5) (3.3,1.4) (4.1,2.45)};
  \draw (1.7,0) -- (1.7,-0.09);
  \node[below] at (1.7,-0.36) {$\tfrac12$};
  \node[below left] at (0.05,-0.1) {$0$};
  \node[primary] at (3.95,2.62) {$f_1$};
  \node[danger] at (3.5,0.45) {$f_2$};
  % правая панель: ét(F)
  \begin{scope}[xshift=6.4cm]
    \node[anchor=south west] at (-0.3,3.05) {$\et(\Shv{F})$};
    \draw[line width=1.4pt] (-0.3,2.1) -- (1.7,2.1);
    \draw[fig leafA,smooth] plot coordinates
      {(1.7,2.1) (2.7,2.45) (4.3,2.85)};
    \draw[fig sect,smooth] plot coordinates
      {(1.7,2.1) (2.7,1.65) (4.3,1.25)};
    \fill (2.7,2.45) circle (2.6pt);
    \fill (2.7,1.65) circle (2.6pt);
    \draw[<->,gray!70!black,line width=0.6pt] (2.7,1.82) -- (2.7,2.28);
    \node[fig note,align=center] at (3.35,0.55)
      {окрестности\\пересекаются};
    \node[primary] at (4.0,2.45) {$f_1$};
    \node[danger] at (3.7,1.25) {$f_2$};
    \draw[fig fiber] (1.7,2.1) -- (1.7,0.15);
    \node at (1.7,-0.36) {$\tfrac12$};
  \end{scope}
\end{tikzpicture}
\end{risunok}

Пояснение к рисунку~\ref{fig:nonhausdorff}. Слева --- графики
двух функций: синий $f_1$ и красный пунктирный $f_2$ совпадают
на отрезке $[0,\tfrac12]$ (участок одного цвета) и дальше
расходятся. Справа --- как это выглядит в $\et(\Shv{F})$: слева
от отметки $\tfrac12$ обе кривые лежат в одном и том же листе
(общий чёрный отрезок), а справа расходятся на два разных листа.
Над точкой $\tfrac12$ два ростка --- две разные точки (чёрные
точки, соединённые стрелкой), но любые их окрестности
обязательно пересекаются, раз листы слева соединены. Отсюда:
$\et(\Shv{F})$ не хаусдорфово.

\begin{bookbox}
\booksrc{Годеман, гл.~II, \S\,1, п.~1.2, с.~130---133: топология
  вводится как <<наиболее тонкая, при которой отображения $s$
 непрерывны>>, и доказывается критерий открытости.
 \textbf{Теорема 1.2.1 (Годеман, с.~132):} всякий пучок множеств с
 базой $X$ изоморфен пучку ростков некоторого накрывающего пространства
 над $X$, определённого однозначно с точностью до изоморфизма.
\par\smallskip
 Манин, \S\,2.1.8, с.~113---114 --- эквивалентное определение пучка:
 пара $(Y_{\Shv{F}},r\colon Y_{\Shv{F}}\to X)$, где $r$ локально ---
 гомеоморфизм, а сечения $\Shv{F}(U)$ --- отображения $s\colon U\to
 Y_{\Shv{F}}$ с $r\circ s=\id$. Обратно, из пучка получается
 $Y_{\Shv{F}}=\coprod_x\Shv{F}_x$ с топологией, порождённой сечениями.
}
\end{bookbox}

%==================================================================
\section[Пучок сечений $\Shv{F}^{\#}$]
 {Пучок сечений $\Shv{F}^{\#}$}
%==================================================================

\begin{defn}[пучок $\Shv{F}^{\#}$]
\begin{equation}\label{eq:Fsharp}
 \Shv{F}^{\#}(U):=\bigl\{\,s\colon U\to\et(\Shv{F})\ \bigm|\ s
 \text{ непрерывно},\ \pi\circ s=\id_U\,\bigr\}.
\end{equation}
Отображения ограничения --- обычные ограничения функций: при
$V\subset U$ отображение $\Shv{F}^{\#}(U)\to\Shv{F}^{\#}(V)$ ---
$s\mapsto s|_{V}$.
\end{defn}

Разберём два условия в~\eqref{eq:Fsharp}:
\begin{conditions}
\item $\pi\circ s=\id_U$ означает, что $s(x)$ лежит над $x$, т.\,е.
 $s(x)\in\Shv{F}_x$: сечение выбирает росток в каждой точке.
\item Непрерывность означает (по теореме~\ref{thm:open}), что для любого
 открытого $G\subset\et(\Shv{F})$ множество $\{x\in U: s(x)\in G\}$
 открыто. Это и есть <<плавность>> выбора.
\end{conditions}

В обеих записях встречается уточнение <<$\pi\circ s=\id_U$», а не
<<$(\pi\circ s)(U)=\id_U$»: имеется в виду тождественность как
отображений, т.\,е. равенство в каждой точке.

\begin{rem}[сравнение с $\Shv{F}(U)$]
Для предпучка определено каноническое отображение
\begin{equation}\label{eq:i}
 i\colon\Shv{F}(U)\longrightarrow\Shv{F}^{\#}(U),\qquad
 i(s)=\tilde s,\quad \tilde s(x)=s_x.
\end{equation}
Предпучок <<теряет>> часть сечений (отождествляет те, что локально
совпадают), а пучок <<добавляет>> (склеивает согласованные локальные
наборы). Между этим и стоит весь разговор о прокачивании.
\end{rem}

\begin{bookbox}
\booksrc{Годеман, Замечание 1.2.3, с.~133: каждый предпучок
 $\Shv{F}$ канонически определяет пучок $\tilde{\Shv{F}}$,
 \emph{порождённый предпучком $\Shv{F}$}, и канонический гомоморфизм
 предпучков $\Shv{F}\to\tilde{\Shv{F}}$ --- это и есть наше
 $i\colon\Shv{F}\to\Shv{F}^{\#}$. Годеман, гл.~II, \S\,1, п.~1.4,
 с.~134: <<простые пучки>> --- частный случай описанной конструкции.
}
\end{bookbox}

%==================================================================
\section{Проверка аксиом (F1) и (F2) для $\Shv{F}^{\#}$}\label{sec:axioms}
%==================================================================

Пока не доказано, что $\Shv{F}^{\#}$ --- пучок, все рассуждения о нём
--- лишь предположение. Проверим аксиомы.

\begin{rem}[алгебраическая структура]
На $\Shv{F}^{\#}(U)$ вводится поэлементная абелева структура (если
$\Shv{F}$ --- предпучок групп):
\[
 (s_1+s_2)(x)=s_1(x)+s_2(x),\qquad
 0(x)=0_{\Shv{F}_x}.
\]
Нулевой элемент --- непрерывное сечение, постоянно равное нулевому
ростку; оно непрерывно, потому что прообраз любого открытого ---
либо $\varnothing$, либо всё $U$. Остаётся проверить (F1) и (F2).
\end{rem}

\subsection{Аксиома (F1): единственность}
\textbf{Задача.} $U=\bigcup_i U_i$, $s',s''\in\Shv{F}^{\#}(U)$ и
$s'|_{U_i}=s''|_{U_i}$ для всех $i$. Доказать $s'=s''$.

\textbf{Доказательство.} Обозначим
\[
 E:=\{\,x\in U:\ s'(x)=s''(x)\,\}.
\]
Шаг 1: $E$ открыто. Действительно, пусть $s'(x)=s''(x)=\xi$. Поскольку
$\pi$ --- локальный гомеоморфизм (теорема~\ref{thm:localhomeo}), найдётся
открытая окрестность $N\ni\xi$ такая, что $\pi|_{N}$ --- гомеоморфизм
на открытое $V\ni x$. По непрерывности $s'$ и $s''$ множества
$(s')^{-1}(N)$ и $(s'')^{-1}(N)$ открыты и содержат $x$; их пересечение
$W$ открыто, и для любого $y\in W$ оба ростка $s'(y),s''(y)$ лежат в
$N$ и проецируются в $y$. Так как $\pi|_N$ инъективно,
$s'(y)=s''(y)$. Значит $W\subset E$, т.\,е. $E$ открыт.

Шаг 2: по условию $U_i\subset E$ для каждого $i$, а $U=\bigcup_iU_i$,
поэтому $U\subset E$, т.\,е. $s'=s''$ на всём $U$.

Заметьте, что <<закрытость $E$» здесь не нужна: достаточно открытости
$E$ и того, что $U_i\subset E$.
\hfill$\square$

\subsection{Аксиома (F2): склейка}
\textbf{Задача.} $U=\bigcup_iU_i$ и даны
$s_i\in\Shv{F}^{\#}(U_i)$, согласованные:
$s_i|_{U_i\cap U_j}=s_j|_{U_i\cap U_j}$.

\textbf{Доказательство.} Определим $s\colon U\to\et(\Shv{F})$ по правилу
$s(x):=s_i(x)$ для любого $i$ с $x\in U_i$.
\begin{conditions}
\item \emph{Согласованность выбора}: если $x\in U_i\cap U_j$, то
 $s_i(x)=s_j(x)$, потому что сечения $s_i$ и $s_j$ совпадают на
 $U_i\cap U_j$ (это их ограничение одно и то же, а сечение
 определяется значениями). Значит $s$ --- функция.
\item \emph{$\pi\circ s=\id$}: выполняется по построению, ведь
 $s_i(x)\in\Shv{F}_x$.
\item \emph{Непрерывность}: пусть $G\subset\et(\Shv{F})$ открыто. Тогда
 \[
   s^{-1}(G)\cap U_i=s_i^{-1}(G)
 \]
 открыто в $U_i$ (непрерывность $s_i$), а $U_i$ открыто в $U$,
 значит $s^{-1}(G)\cap U_i$ открыто и в $U$. Наконец
 \[
   s^{-1}(G)=\bigcup_i\bigl(s^{-1}(G)\cap U_i\bigr)
 \]
 --- объединение открытых в $U$, значит открыто.
\item \emph{Ограничения}: по построению $s|_{U_i}=s_i$.
\end{conditions}
Уникальность $s$ тривиальна: сечение определяется своими значениями.
\hfill$\square$

Итак, $\Shv{F}^{\#}$ --- действительно пучок (и пучок абелевых групп,
если $\Shv{F}$ --- предпучок групп).

\begin{risunok}{Проверка аксиом для $\Shv{F}^{\#}$: слева --- (F1)
 (равны на каждом $U_i$ $\Rightarrow$ равны), справа --- (F2)
 (согласованные $s_1,s_2$ склеиваются в $s$). \label{fig:axioms}}
\resizebox{\linewidth}{!}{\begin{tikzpicture}[>=Stealth,font=\small]
  % левая панель: (F1)
  \begin{scope}
    \draw[fig base] (0,0) ellipse (1.9 and 0.6);
    \begin{scope}
      \clip (0,0) ellipse (1.9 and 0.6);
      \fill[primary!10] (-0.7,0) ellipse (1.35 and 0.55);
      \fill[success!10] (0.7,0) ellipse (1.35 and 0.55);
    \end{scope}
    \draw[fig open] (-0.7,0) ellipse (1.35 and 0.55);
    \draw[fig open2] (0.7,0) ellipse (1.35 and 0.55);
    \node[primary] at (-1.7,0.8) {$U_1$};
    \node[success] at (1.7,0.8) {$U_2$};
    \draw[fig fiber] (-1.0,0.51) -- (-1.0,2.35);
    \draw[fig fiber] ( 1.0,0.51) -- ( 1.0,2.35);
    \fill (-1.0,1.75) circle (2.6pt);
    \fill ( 1.0,1.95) circle (2.6pt);
    \draw[fig sect,smooth] plot coordinates
      {(-2.1,1.4) (-1.0,1.75) (0,1.5) (1.0,1.95) (2.1,1.7)};
    \draw[danger,line width=1pt,dashed,smooth] plot coordinates
      {(-2.1,1.75) (-1.0,1.75) (0,2.05) (1.0,1.95) (2.1,2.15)};
    \node[danger] at (-2.45,1.2) {$s'$};
    \node[danger,font=\small] at (2.45,2.3) {$s''$};
    \node[font=\footnotesize] at (0,-1.05)
      {(F1): равны на каждом $U_i$ $\Rightarrow$ равны};
  \end{scope}
  % правая панель: (F2)
  \begin{scope}[xshift=6.6cm]
    \draw[fig base] (0,0) ellipse (2.4 and 0.7);
    \begin{scope}
      \clip (0,0) ellipse (2.4 and 0.7);
      \fill[primary!10] (-0.85,0) ellipse (1.55 and 0.65);
      \fill[success!10] (0.85,0) ellipse (1.55 and 0.65);
    \end{scope}
    \draw[fig open] (-0.85,0) ellipse (1.55 and 0.65);
    \draw[fig open2] (0.85,0) ellipse (1.55 and 0.65);
    \node[primary] at (-2.2,0.85) {$U_1$};
    \node[success] at (2.2,0.85) {$U_2$};
    \draw[fig fiber] (-1.8,0.463) -- (-1.8,2.65);
    \draw[fig fiber] ( 0,0.7)     -- ( 0,2.65);
    \draw[fig fiber] ( 1.8,0.463) -- ( 1.8,2.65);
    \fill (-1.8,1.45) circle (2.6pt);
    \fill ( 0,2.15)   circle (2.6pt);
    \fill ( 1.8,1.65) circle (2.6pt);
    \draw[fig leafA,smooth] plot coordinates
      {(-2.4,1.3) (-1.8,1.45) (-0.9,1.85) (0,2.15)};
    \draw[fig leafB,smooth] plot coordinates
      {(0,2.15) (0.9,2.0) (1.8,1.65) (2.4,1.5)};
    \node[primary] at (-1.3,1.22) {$s_1$};
    \node[success] at (1.25,1.28) {$s_2$};
    \node[danger] at (0.42,2.52) {$s$};
    \node[font=\footnotesize] at (0,-1.15)
      {(F2): $s_1,s_2$ согласованы $\Rightarrow$ есть $s$};
  \end{scope}
\end{tikzpicture}}
\end{risunok}

Пояснение к рисунку~\ref{fig:axioms}. Слева --- аксиома (F1):
две кривые $s'$ (сплошная) и $s''$ (пунктирная) заранее разные,
но над каждым $U_i$ их значения совпадают --- это и есть общие
точки на волокнах; (F1) утверждает, что тогда $s'=s''$
отождествляются полностью. Справа --- аксиома (F2): локальные
сечения $s_1$ (синее, над $U_1$) и $s_2$ (зелёное, над $U_2$)
согласованы в общей части $U_1\cap U_2$ --- обе кривые проходят
через одну и ту же точку на общем волокне, --- и их склейка
даёт глобальное сечение $s$.

\sectionsrc{photo\_2 (стр.~10): <<внезапная вставка с F1 и F2
 для $\et(\Shv{F})$» --- с рисунками, на которых листы $\Shv{F}^{\#}$
 над покрытием $U=U_i\cup U_j$ показаны <<зелёными и красными»;
 отмечено, что мало проверить равенство образов, надо именно равенство
 сечений.}

\begin{bookbox}
\booksrc{Годеман, п.~3.9, с.~183---184: каноническое отображение
 $\Shv{F}(X)\to\Shv{F}^{\#}(X)$ инъективно, каково бы ни было
 $\Shv{F}$, если выполняется (F1); если для $\Shv{F}$ выполнена (F2)
 и $X$ паракомпактно, то это отображение --- на всё множество.
 Общая идея: в $\Shv{F}(X)$ <<локально равные» сечения отождествляются,
 а в $\Shv{F}^{\#}(X)$ <<добавляются» сечения, склеенные из
 согласованных локальных наборов.
}
\end{bookbox}

%==================================================================
\section{Пример: постоянный предпучок}\label{sec:const}
%==================================================================

\begin{exm}[простой пучок]\label{exm:const}
Пусть $A$ --- абелева группа, а предпучок $\Shv{F}$ определён формулой
$\Shv{F}(U)=A$ для всех непустых открытых $U$ (и
$\Shv{F}(\varnothing)=0$), все ограничения $\rho$ --- тождественные.
Тогда:
\begin{bullets}
\item $\Shv{F}_x=A$ для всякой $x$ (предел по окрестностям ничего не
 меняет, ведь все $\rho$ --- тождества), поэтому
 $\et(\Shv{F})=X\times A$, и $\pi\colon X\times A\to X$ --- обычная
 проекция;
\item топология на $X\times A$ --- произведение топологии $X$ и
  дискретной топологии на $A$ (это как раз наиболее тонкая топология,
 при которой все <<константные>> сечения непрерывны);
\item $\Shv{F}^{\#}(U)$ --- непрерывные отображения $U\to A$, т.\,е.
 \emph{локально постоянные} функции, а на связном $U$ --- просто
 константы.
\end{bullets}
Проверка (F1) для этого примера --- обычное утверждение из топологии:
локально постоянные функции на связном множестве постоянны, а на
произвольном --- определяются значениями на связных компонентах.
\end{exm}

\begin{risunok}{Постоянный предпучок: $\et(\Shv{F})=X\times A$ с
 дискретным $A$; сечение --- локально постоянная функция $U\to A$.
  \label{fig:const}}
\begin{tikzpicture}[>=Stealth,font=\small]
  % база X
  \draw[fig base] (0,0) ellipse (3.2 and 0.7);
  \node[below] at (0,-0.85) {$X$};
  % уровни {x} x {a}
  \foreach \y in {1.5,2.5,3.5} \draw[fig fiber] (-3.2,\y) -- (3.2,\y);
  \foreach \y/\a in {1.5/{a_1},2.5/{a_2},3.5/{a_3}} {
    \node[anchor=east,font=\footnotesize] at (-3.35,\y)
      {$\{x\}\times\{\a\}$};
    \foreach \x in {-2.6,-1.6,-0.6,0.4,1.4,2.4}
      \fill (\x,\y) circle (2.2pt);
  }
  % волокна
  \foreach \x/\top in {-2.6/0.408,-1.6/0.606,-0.6/0.687,
                       0.4/0.694,1.4/0.629,2.4/0.463}
    \draw[fig fiber] (\x,\top) -- (\x,3.5);
  % открытые U_1, U_2
  \begin{scope}
    \clip (0,0) ellipse (3.2 and 0.7);
    \fill[primary!10] (-2.1,0) ellipse (1.45 and 0.65);
    \fill[success!10] (1.5,0) ellipse (1.45 and 0.65);
  \end{scope}
  \draw[fig open] (-2.1,0) ellipse (1.45 and 0.65);
  \draw[fig open2] (1.5,0) ellipse (1.45 and 0.65);
  \node[primary] at (-2.1,-1.1) {$U_1$};
  \node[success] at (1.5,-1.1) {$U_2$};
  % локально постоянные сечения
  \draw[fig sect] (-3.15,2.5) -- (-0.7,2.5);
  \draw[fig sect] (0.1,1.5) -- (2.9,1.5);
  \node[danger] at (-1.9,2.68) {$a_2$};
  \node[danger] at (1.5,1.68) {$a_1$};
  % проекция
  \draw[->,thick] (3.75,3.5) -- (3.75,0.5);
  \node at (3.95,2.0) {$\pi$};
  \node[anchor=south east] at (-3.3,3.85) {$X\times A=\et(\Shv{F})$};
\end{tikzpicture}
\end{risunok}

Пояснение к рисунку~\ref{fig:const}. Овал внизу --- база $X$;
всё вместе --- $\et(\Shv{F})=X\times A$, где $A$ дискретно,
поэтому <<этажёрка>> состоит из отдельных горизонтальных
рядов точек:
горизонтальные пунктирные линии --- уровни $\{x\}\times\{a\}$
для трёх значений $a_1,a_2,a_3$. Синее и зелёное --- открытые
$U_1,U_2\subset X$. Красные отрезки --- сечение: над $U_1$ оно
постоянно равно $a_2$ (лежит на одном уровне), над $U_2$ ---
равно $a_1$: именно локально постоянная функция $U\to A$.
Стрелка $\pi$ --- проекция на базу.

\begin{bookbox}
\booksrc{Годеман, гл.~II, \S\,1, п.~1.4 <<Простые пучки>>, с.~134:
 простой пучок с базой $X$ и слоем $A$ --- пучок, порождённый предпучком
 $U\mapsto A$ с тождественными ограничениями;
 $\et=X\times A$ --- произведение топологии $X$ и дискретной топологии
 на $A$; сечения над $U$ канонически отождествляются с локально
 постоянными отображениями $U\to A$.
}
\end{bookbox}

%==================================================================
\section[Основное утверждение: $\Hom(\Shv{F},\Shv{G})\simeq\Hom(\Shv{F}^{\#},\Shv{G})$]
 {Основное утверждение: $\Hom(\Shv{F},\Shv{G})\simeq
  \Hom(\Shv{F}^{\#},\Shv{G})$}
%==================================================================

\begin{thm}[главное утверждение лекции]\label{thm:hom}
Пусть $\Shv{F}$ --- предпучок на $X$, а $\Shv{G}$ --- \emph{пучок} на
$X$. Тогда отображение
\begin{equation}\label{eq:Phi}
 \Phi\colon\Hom(\Shv{F}^{\#},\Shv{G})\longrightarrow
 \Hom(\Shv{F},\Shv{G}),\qquad \varphi\longmapsto\varphi\circ i,
\end{equation}
где $i\colon\Shv{F}\to\Shv{F}^{\#}$ --- каноническое, является
изоморфизмом. Более того, изоморфизм естествен: коммутативны
квадраты диаграмм~\ref{dgm:natG} и~\ref{dgm:natF}.
\end{thm}

Заметим направление: изоморфизм идёт \emph{из} морфизмов в пучок
$\Shv{F}^{\#}$ \emph{в} морфизмы в произвольный предпучок. Гомоморфизм
пучка определяется своими значениями на сечениях, а сечений
$\Shv{F}^{\#}(U)$ <<достаточно много>>: именно это и доказывается.

\subsection{Шаг 1: инъективность $\Phi$}
Для $\varphi\colon\Shv{F}^{\#}\to\Shv{G}$ положим
$\varphi':=\varphi\circ i\colon\Shv{F}\to\Shv{G}$ --- диаграмма
\ref{dgm:univ}.

\begin{diagram}{Диаграмма универсальности: $\varphi'=\varphi\circ i$
 --- сжатие $\varphi$ по каноническому $i$.\label{dgm:univ}}
\begin{tikzcd}[cd style,column sep=huge,row sep=huge]
 \Shv{F} \ar[r,"i"] \ar[dr,"\varphi'"'] & \Shv{F}^{\#} \ar[d,"\varphi"'] \\
 & \Shv{G}
\end{tikzcd}
\end{diagram}

Пояснение к диаграмме~\ref{dgm:univ}. Треугольник читается как
<<сначала сжать по $i$, затем применить $\varphi$>>: диагональ
$\varphi'=\varphi\circ i$ --- это и есть $\Phi(\varphi)$, переход
морфизма пучка в морфизм предпучков. Коммутативность треугольника
--- просто определение $\varphi'$; всё содержание утверждения
в том, что $\varphi'$ корректно определён и что из
$\varphi\circ i=\psi\circ i$ следует $\varphi=\psi$ ---
доказывается сразу ниже.

Если $\varphi\circ i=\psi\circ i$, то $\varphi=\psi$: действительное
значение $\varphi$ на произвольном сечении $\tilde s\in\Shv{F}^{\#}(U)$
уже определено, потому что $\tilde s$ локально является образом
некоторого $s\in\Shv{F}(V)$, а $\varphi$ --- гомоморфизм \emph{пучка}
$\Shv{G}$, поэтому его поведение локально, а значит и везде,
определяется значениями $i(s)$. Разделение точек $\Phi(\varphi)$ ---
это инъективность.

\subsection{Шаг 2: сюръективность $\Phi$}
\subsubsection*{Построение $\beta$}

Пусть дан $\alpha\colon\Shv{F}\to\Shv{G}$ --- гомоморфизм предпучков.
Надо построить $\beta\colon\Shv{F}^{\#}\to\Shv{G}$ с
$\beta\circ i=\alpha$.

Начнём с <<наивной>> формулы. Для $U\subset X$ и
$s\in\Shv{F}(U)$ положим
\begin{equation}\label{eq:beta-naive}
 \beta\bigl(i(s)\bigr):=\alpha_U(s)\in\Shv{G}(U).
\end{equation}
Это хорошо согласуется с ограничениями, потому что $\alpha$ ---
гомоморфизм предпучков: $\beta(i(s))|_{V}=\alpha_V(s|_V)$.

\subsubsection*{Проблема единственности}

Возникает вопрос: а что, если $s_1\neq s_2$, но $i(s_1)=i(s_2)$? Тогда
по~\eqref{eq:beta-naive} у $\beta$ получится два разных значения
$g_1=\alpha(s_1)$ и $g_2=\alpha(s_2)$, и отображение не будет
определено. Схема на листе~11: $s_1\to\tilde s\leftarrow s_2$.
Покажем, что в таком случае $g_1=g_2$.

\begin{lem}[локальное равенство]\label{lem:local}
Пусть $s,t\in\Shv{F}(U)$ и $i(s)=i(t)$ в $\Shv{F}^{\#}(U)$. Тогда
$\alpha_U(s)=\alpha_U(t)$ в $\Shv{G}(U)$.
\end{lem}

\begin{proof}
По условию ростки совпадают: $s_x=t_x$ для всех $x\in U$. По
определению прямого предела~\eqref{eq:stalk} это означает: для каждого
$x\in U$ существует открытая окрестность $V_x\subset U$ такая, что
\begin{equation}\label{eq:same}
 \res^{V_x}_{U}(s)=\res^{V_x}_{U}(t),
 \qquad\text{т.\,е.}\qquad s|_{V_x}=t|_{V_x}.
\end{equation}
Семейство $\{V_x\}_{x\in U}$ покрывает $U$. Применив к~\eqref{eq:same}
гомоморфизм предпучка $\alpha$ (точнее, его компоненту на $V_x$), получим
\[
 \alpha_{V_x}\bigl(s|_{V_x}\bigr)=\alpha_{V_x}\bigl(t|_{V_x}\bigr),
\]
потому что $\alpha$ переводит равные элементы в равные. С другой
стороны, коммутативность диаграммы ограничений для $\alpha$
(диаграмма~\ref{dgm:natur}) даёт
\[
 \alpha_{V_x}\bigl(s|_{V_x}\bigr)=\bigl(\alpha_U(s)\bigr)|_{V_x},
 \qquad
 \alpha_{V_x}\bigl(t|_{V_x}\bigr)=\bigl(\alpha_U(t)\bigr)|_{V_x}.
\]
Значит $\alpha_U(s)$ и $\alpha_U(t)$ совпадают на каждом $V_x$:
$\bigl(\alpha_U(s)\bigr)|_{V_x}=\bigl(\alpha_U(t)\bigr)|_{V_x}$.
Осталось применить аксиому \emph{(F1) пучка $\Shv{G}$} к покрытию
$U=\bigcup_xV_x$: два глобальных сечения, совпадающие на каждом куске
покрытия, равны. Итак, $\alpha_U(s)=\alpha_U(t)$.
\end{proof}

\begin{diagram}{Натурность гомоморфизма $\alpha\colon\Shv{F}\to\Shv{G}$:
 квадраты коммутативны для любой пары $W\subset W'$, в частности для
 $V_x\subset U$. Слева направо --- ограничения, сверху вниз ---
 компоненты $\alpha$.\label{dgm:natur}}
\begin{tikzcd}[cd style]
 \Shv{F}(U) \ar[r,"\res^{V_x}_{U}"] \ar[d,"\alpha_U"'] &
 \Shv{F}(V_x) \ar[d,"\alpha_{V_x}"] \\
 \Shv{G}(U) \ar[r,"\res^{V_x}_{U}"'] & \Shv{G}(V_x)
\end{tikzcd}
\end{diagram}

Пояснение к диаграмме~\ref{dgm:natur}. Это квадрат натуральности
гомоморфизма предпучков $\alpha$: два маршрута из $\Shv{F}(U)$
в $\Shv{G}(V_x)$ --- «сначала ограничить, потом применить
$\alpha$» (вверх-вправо-вниз) и «сначала применить $\alpha$ на
$U$, потом ограничить» (влево-вниз-направо) --- дают одно и то
же. Это ровно условие $\alpha_U(s)|_{V_x}=\alpha_{V_x}(s|_{V_x})$;
в доказательстве квадрат применяется к паре $V_x\subset U$.

Лемма~\ref{lem:local} означает: формула~\eqref{eq:beta-naive} не зависит
от того, каким из $s$ с $\tilde s=i(s)$ мы воспользуемся.

\subsubsection*{Определение $\beta$ на произвольном сечении}

Теперь возьмём произвольное $\tilde s\in\Shv{F}^{\#}(U)$, а не только
вида $i(s)$. Почему оно локально такого вида? Потому что непрерывность:
для точки $x\in U$ росток $\tilde s(x)\in\Shv{F}_x$ --- это класс
некоторой пары $(V,s)$, $s\in\Shv{F}(V)$; множество точек, в которых
$\tilde s$ совпадает с $i(s)$, открыто (см.~шаг~1 доказательства
(F1) в \S\,\ref{sec:axioms}---там показано, что множество точек
совпадения двух непрерывных сечений открыто). Эти открытые множества
покрывают $U$, значит $U=\bigcup_\alpha W_\alpha$ и
$\tilde s|_{W_\alpha}=i(s_\alpha)$ для некоторых
$s_\alpha\in\Shv{F}(W_\alpha)$.

Положим
\[
 t_\alpha:=\alpha_{W_\alpha}(s_\alpha)\in\Shv{G}(W_\alpha).
\]
На пересечении $W_\alpha\cap W_\beta$ выполняется
$i(s_\alpha)|=i(s_\beta)|$, т.\,е. локальное равенство сечений
$\Shv{F}$; по лемме~\ref{lem:local} (вернее, по её локальной версии)
$\alpha(s_\alpha)$ и $\alpha(s_\beta)$ совпадают на
$W_\alpha\cap W_\beta$. Итак, семейство $(t_\alpha)$ --- согласованное,
и аксиома \emph{(F2) пучка $\Shv{G}$} даёт единственное
$t\in\Shv{G}(U)$ с $t|_{W_\alpha}=t_\alpha$. Полагаем
\begin{equation}\label{eq:beta}
 \beta(\tilde s):=t.
\end{equation}

Независимость от выбора покрытия: если есть другое покрытие
$U=\bigcup_\beta W'_\beta$, то на пересечениях $W_\alpha\cap W'_\beta$
оба определения совпадают (там $\tilde s$ --- образ одного и того же
сечения), и снова применяется (F2), либо (F1) --- глобальное сечение
$\Shv{G}(U)$ определяется по значениям на покрытии однозначно.

Итого: $\beta\colon\Shv{F}^{\#}\to\Shv{G}$ --- гомоморфизм пучков,
и $\beta\circ i=\alpha$ по построению~\eqref{eq:beta}. Значит
$\Phi(\beta)=\alpha$, т.\,е. $\Phi$ --- сюръективно.

\subsubsection*{Схема склейки}

Глобальное сечение $t$ из~\eqref{eq:beta} собирается из локальных
$t_\alpha$ ровно так, как показано на диаграмме~\ref{dgm:glue}: значения
на $W_\alpha$ и $W_\beta$ дают одно и то же на пересечении, и (F2)
<<склеивает>> их в $t\in\Shv{G}(U)$.

\begin{diagram}{Склейка: из локальных сечений $t_x\in\Shv{G}(V_x)$,
 согласованных на $V_x\cap V_y$, аксиома (F2) пучка $\Shv{G}$ даёт
 единственное $t\in\Shv{G}(U)$.\label{dgm:glue}}
\begin{tikzcd}[cd style]
 & \Shv{G}(U) \ar[dl,"\res^{V_x}_{U}"'] \ar[dr,"\res^{V_y}_{U}"] & \\
 \Shv{G}(V_x) \ar[dr,"\res"'] & & \Shv{G}(V_y) \ar[dl,"\res"] \\
  & \Shv{G}(V_x\cap V_y) &
\end{tikzcd}
\end{diagram}

Пояснение к диаграмме~\ref{dgm:glue}. Ромб показывает действие
аксиомы (F2) в обе стороны. Вниз: глобальное сечение
$t\in\Shv{G}(U)$ ограничивается до $V_x$ и $V_y$, и оба
полученных локальных значения совпадают на пересечении
$V_x\cap V_y$ (две нижние стрелки сходятся к одному элементу).
Вверх --- сама склейка: если заданы согласованные локальные
значения $t_x$ и $t_y$, то (F2) даёт единственное
$t\in\Shv{G}(U)$, из которого они получаются ограничением.

\sectionsrc{photo\_3---photo\_4 (стр.~11---12):
 проблема $s_1\neq s_2$, $i(s_1)=i(s_2)$; выбор $V_x$ <<по определению
 прямого предела>>; натурная диаграмма; вывод
 <<$\forall x,y\in U\ \exists V_x,V_y$ и $g_x\in\Shv{G}(V_x)$,
 $g_y\in\Shv{G}(V_y)$ с $g_x|_{V_x\cap V_y}=g_y|_{V_x\cap V_y}$;
 применяем (F2) к пучку $\Shv{G}$ и получаем
 $g\in\Shv{G}(U)$, определённый однозначно $\Rightarrow$ отображение
 $\beta$ задано корректно».}

\subsection{Шаг 3: естественность}
Для гомоморфизма пучков $f\colon\Shv{G}_1\to\Shv{G}_2$ и для
гомоморфизма предпучков $g\colon\Shv{F}_1\to\Shv{F}_2$ коммутативны
квадраты диаграмм~\ref{dgm:natG} и~\ref{dgm:natF}.

\begin{diagram}{Естественность $\Phi$ по $\Shv{G}$: сжатие по $i$
 совместно с композицией с $f$.\label{dgm:natG}}
\begin{tikzcd}[cd style]
 \Hom(\Shv{F},\Shv{G}_1) \ar[r,"f\circ-"] \ar[d,"\Phi"'] &
 \Hom(\Shv{F},\Shv{G}_2) \ar[d,"\Phi"] \\
 \Hom(\Shv{F}^{\#},\Shv{G}_1) \ar[r,"f\circ-"'] &
 \Hom(\Shv{F}^{\#},\Shv{G}_2)
\end{tikzcd}
\end{diagram}

Пояснение к диаграмме~\ref{dgm:natG}. Квадрат естественности
$\Phi$ по пучку: верхняя стрелка --- домножение морфизма
предпучков на $f\colon\Shv{G}_1\to\Shv{G}_2$, нижняя --- то же
для морфизмов в $\Shv{G}_2$, вертикали --- <<сжатие>> $\Phi$
по каноническому $i$. Коммутативность означает:
$f\circ(\varphi\circ i)=(f\circ\varphi)\circ i$, т.\,е.\ неважно,
сжимать до или после домножения на $f$.

\begin{diagram}{Естественность $\Phi$ по $\Shv{F}$: сжатие по $i_2$
 совместно с композиции с $g$.\label{dgm:natF}}
\begin{tikzcd}[cd style]
 \Hom(\Shv{F}_1,\Shv{G}) \ar[r,"-\circ i_2"] \ar[d,"\Phi"'] &
 \Hom(\Shv{F}_2,\Shv{G}) \ar[d,"\Phi"] \\
 \Hom(\Shv{F}_1^{\#},\Shv{G}) \ar[r,"-\circ i_2"'] &
 \Hom(\Shv{F}_2^{\#},\Shv{G})
\end{tikzcd}
\end{diagram}

Пояснение к диаграмме~\ref{dgm:natF}. Зеркальная ситуация:
стрелки $-\circ i_2$ домножают морфизм предпучков
$g\colon\Shv{F}_1\to\Shv{F}_2$, вертикали --- снова $\Phi$.
Коммутативность --- это $(\varphi\circ g)\circ i_1
=(\varphi\circ i_2)\circ g$: сжимать сначала $g$, потом $\varphi$,
или наоборот, --- одинаково. Вместе две диаграммы
(\ref{dgm:natG} и \ref{dgm:natF}) и означают, что изоморфизм
$\Phi$ естественен по обоим аргументам.

Здесь $i_2\colon\Shv{F}_2\to\Shv{F}_2^{\#}$ --- каноническое
отображение, а стрелки $-\circ i_2$ индуцированы гомоморфизмом
предпучков $g$. Проверка прямая: $\Phi(f\circ\varphi)=
f\circ\varphi\circ i=f\circ\Phi(\varphi)$ и
$\Phi(\varphi\circ g)=(\varphi\circ g)\circ i_1=(\varphi\circ i_2)\circ
g=\Phi(\varphi)\circ g$.

\begin{rem}[что это значит категорно]
Если задан функтор $\PSh(X)\to\Sh(X)$, отсылающий предпучок к пучку,
порождённому им, то изоморфизм теоремы~\ref{thm:hom} ---
это \emph{свойство универсальности} этого функтора: включение
$\Sh(X)\hookrightarrow\PSh(X)$ имеет левый сопряжённый
$\Shv{F}\mapsto\Shv{F}^{\#}$. Коммутативность диаграмм
\ref{dgm:natG}, \ref{dgm:natF} --- naturality этого построения;
именно она позволяет говорить, что <<прокачивание>> задано
\emph{функторно}, а не произвольным выбором.
\end{rem}

\begin{bookbox}
\booksrc{Годеман, гл.~I, п.~1.7, с.~24---25: категории, ковариантные и
 контравариантные функторы, естественное преобразование (система
 $T(X)$, которая делает диаграмму с $F(f),G(f)$ коммутативной); пример:
 морфизм $t\colon A\to B$ индуцирует морфизм функторов
 $\Hom(B,X)\to\Hom(A,X)$ передачей $u\mapsto u\circ t$.
 Годеман, гл.~I, п.~1.9, с.~30---31: предпучок --- контравариантный
 функтор категории открытых множеств; категория предпучков абелевых
 групп --- абелева категория, функтор $\Shv{F}\mapsto\Shv{F}(U)$ точен,
 поэтому точность последовательности предпучков эквивалентна
 поэлементной точности.
\par\smallskip
 Манин, \S\,1.16.5---1.16.6, с.~103: предпучок как контравариантный
 функтор $\mathrm{Top}_E\to\mathrm{Sets}$; определение функторного
 морфизма с коммутативной диаграммой; категория
 $\mathrm{Funct}(C,D)$.
}
\end{bookbox}

%==================================================================
\section{<<Факт>> о функторах $\Hom$ и характеризация пучка}
%==================================================================

\subsection{Факт}
\leavevmode\par
\begin{fact}[Факт из лекции]\label{fact:hom}
Пусть $\mathcal{C}$ --- маленькая категория,
$A_1,A_2\in\mathrm{Ob}\,\mathcal{C}$,
$F_1=\Hom(A_1,-)$, $F_2=\Hom(A_2,-)$ --- функторы
$\mathcal{C}\to\mathrm{Sets}$. Если существует натуральное
преобразование $\eta\colon F_1\Rightarrow F_2$ такое, что каждая его
компонент $\eta_X\colon\Hom(A_1,X)\to\Hom(A_2,X)$ является биекцией, то
$A_1\simeq A_2$.
\end{fact}

Иными словами: функтор $\Hom(A,-)$ определяет объект категории
\emph{однозначно}. Интуиция: $A$ --- это <<тот объект, через который
удобно смотреть на все остальные>>, и если два объекта дают одинаковые
<<окна>>, они изоморфны.

\subsection{Доказательство}
\textbf{Шаг 1: достаём $f$ и $g$.}
Рассмотрим компоненту $\eta_{A_1}\colon\Hom(A_1,A_1)\to\Hom(A_2,A_1)$;
она переводит $\id_{A_1}$ в некоторый морфизм
\[
 f\colon A_2\longrightarrow A_1.
\]
Компонента $\eta_{A_2}\colon\Hom(A_1,A_2)\to\Hom(A_2,A_2)$ --- биекция,
поэтому можно рассмотреть обратную $\eta_{A_2}^{-1}$; она переводит
$\id_{A_2}$ в морфизм
\[
 g\colon A_1\longrightarrow A_2.
\]

\textbf{Шаг 2: докажем $g\circ f=\id_{A_2}$.}
Натуральность $\eta$ для морфизма $g\colon A_1\to A_2$ --- это
коммутативность диаграммы~\ref{dgm:fact1}.

\begin{diagram}{Натуральность $\eta$ для $g\colon A_1\to A_2$: стрелки
 $-\circ g$ --- композиция с $g$.\label{dgm:fact1}}
\begin{tikzcd}[cd style]
 \Hom(A_1,A_1) \ar[r,"\eta_{A_1}"] \ar[d,"g\circ-"'] &
 \Hom(A_2,A_1) \ar[d,"g\circ-"] \\
 \Hom(A_1,A_2) \ar[r,"\eta_{A_2}"'] & \Hom(A_2,A_2)
\end{tikzcd}
\end{diagram}

Пояснение к диаграмме~\ref{dgm:fact1}. Это условие натуральности
преобразования $\eta$ при морфизме $g$: маршрут <<сначала
$\eta_{A_1}$, потом скомпоновать с $g$>> (вправо-вниз) равен
маршруту <<сначала скомпоновать с $g$, потом $\eta_{A_2}$>>
(влево-вниз-вправо). Ниже оба маршрута прослеживаются на
элементе $\id_{A_1}$, и равенство даёт $g\circ f=\id_{A_2}$.

Проследим, куда попадает $\id_{A_1}$ по двум маршрутам:
\begin{bullets}
\item сверху направо: $\id_{A_1}\mapsto\eta_{A_1}(\id_{A_1})=f$, затем
 вниз: $f\mapsto g\circ f$;
\item слева вниз, затем направо: $\id_{A_1}\mapsto g$, затем
 $g\mapsto\eta_{A_2}(g)$.
\end{bullets}
Коммутативность даёт $\eta_{A_2}(g)=g\circ f$. С другой стороны, по
построению $g$ мы имели $\eta_{A_2}^{-1}(\id_{A_2})=g$, т.\,е.
$\eta_{A_2}(g)=\id_{A_2}$. Значит
\begin{equation}\label{eq:gf}
 g\circ f=\id_{A_2}.
\end{equation}

\textbf{Шаг 3: строим обратное преобразование $\eta^{-1}$.}
Положим $(\eta^{-1})_X:=\eta_X^{-1}$, т.\,е.
$\Hom(A_2,X)\to\Hom(A_1,X)$. Надо проверить натуральность: для любого
морфизма $\alpha\colon X\to Y$ диаграмма~\ref{dgm:fact2} должна быть
коммутативной.

\begin{diagram}{Натуральность обратного преобразования: компоненты
 $\eta_X^{-1}$ и композиции с $\alpha$.\label{dgm:fact2}}
\begin{tikzcd}[cd style]
 \Hom(A_2,X) \ar[r,"\alpha\circ-"] \ar[d,"\eta_X^{-1}"'] &
 \Hom(A_2,Y) \ar[d,"\eta_Y^{-1}"] \\
 \Hom(A_1,X) \ar[r,"\alpha\circ-"'] & \Hom(A_1,Y)
\end{tikzcd}
\end{diagram}

Пояснение к диаграмме~\ref{dgm:fact2}. Условие натуральности
для обратного преобразования $\eta^{-1}$. Два маршрута из
$\Hom(A_2,X)$ в $\Hom(A_1,Y)$: вправо-вниз --- «сначала
скомпоновать с $\alpha$, потом применить $\eta^{-1}$» ---
и вниз-вправо --- «сначала $\eta^{-1}$, потом композиция
с $\alpha$»; натуральность и есть их равенство. На элементе
$p$ маршруты прослеживаются ниже.

Проверка: пусть $p\colon A_2\to X$. Поскольку $\eta_X$ --- биекция,
$p=\eta_X(q)$ для единственного $q\colon A_1\to X$, т.\,е.
$q=\eta_X^{-1}(p)$. Тогда оба маршрута дают
$\eta_Y^{-1}(\alpha\circ\eta_X(q))$: слева направо --- по определению,
справа налево --- $\eta_Y^{-1}(\alpha\circ p)$. Коммутативность
доказана.

\textbf{Шаг 4: докажем $f\circ g=\id_{A_1}$.}
Рассмотрим диаграмму~\ref{dgm:fact3}, натуральность $\eta^{-1}$ для
морфизма $f\colon A_2\to A_1$, и применим её к $\id_{A_2}$.

\begin{diagram}{Натуральность $\eta^{-1}$ для $f\colon A_2\to A_1$
 (та же диаграмма \ref{dgm:fact2} с частным выбором $X=A_2$, $Y=A_1$;
 стрелки $f\circ-$).\label{dgm:fact3}}
\begin{tikzcd}[cd style]
 \Hom(A_2,A_2) \ar[r,"\eta_{A_2}^{-1}"] \ar[d,"f\circ-"'] &
 \Hom(A_1,A_2) \ar[d,"f\circ-"] \\
 \Hom(A_2,A_1) \ar[r,"\eta_{A_1}^{-1}"'] & \Hom(A_1,A_1)
\end{tikzcd}
\end{diagram}

Пояснение к диаграмме~\ref{dgm:fact3}. Частный случай
диаграммы~\ref{dgm:fact2}: naturality $\eta^{-1}$, применённая
к морфизму $f\colon A_2\to A_1$. Здесь стрелки --- композиция
с $f$; приложение коммутативности к $\id_{A_2}$ даёт
$f\circ g=\id_{A_1}$.

Верхний путь: $\id_{A_2}\mapsto\eta_{A_2}^{-1}(\id_{A_2})=g$, затем
вниз: $g\mapsto f\circ g$. Нижний путь: $\id_{A_2}\mapsto f$, затем
направо: $f\mapsto\eta_{A_1}^{-1}(f)=\eta_{A_1}^{-1}(\eta_{A_1}(\id_{A_1}))
=\id_{A_1}$. Значит
\begin{equation}\label{eq:fg}
 f\circ g=\id_{A_1}.
\end{equation}

Из~\eqref{eq:gf} и~\eqref{eq:fg} следует $A_1\simeq A_2$ (изоморфизм
$f$ с обратным $g$). \hfill$\square$

\subsection{Следствие: как узнать пучок}
\leavevmode\par
\begin{cor}\label{cor:criterion}
Пусть $\Shv{F}$ --- предпучок, $\Shv{F}'$ --- пучок. Если для всех
пучков $\Shv{G}$ изоморфизм $\Hom(\Shv{F},\Shv{G})\simeq\Hom(\Shv{F}',
\Shv{G})$ естествен (т.\,е. коммутативны все диаграммы вида
\ref{dgm:natG} и \ref{dgm:natF}), то $\Shv{F}\simeq\Shv{F}'$ как
предпучки.
\end{cor}

\begin{proof}
Возьмём $\mathcal{C}=\Sh(X)$ и $F_1=\Hom(\Shv{F},-)$,
$F_2=\Hom(\Shv{F}',-)$ как функторы $\Sh(X)\to\mathrm{Sets}$. Данное
естественное преобразование удовлетворяет условиям Факта, поэтому
$\Shv{F}\simeq\Shv{F}'$ в $\mathcal{C}$, т.\,е. как пучки; но $\Shv{F}'$
--- пучок, значит и $\Shv{F}$ изоморфен пучку.
\end{proof}

Это следствие и есть <<принцип единственности>>: два способа
<<прокачать>> предпучок, дающие одинаковые морфизмы во все пучки,
совпадают. В частности, $\Shv{F}^{\#}$ определён однозначно.

\begin{rem}[лемма Ёнеды]
Факт~\ref{fact:hom} --- частный случай знаменитого погружения
категории в категорию представимых функторов (лемма Ёнеды).
\end{rem}

\begin{bookbox}
\booksrc{Манин, \S\,1.16.8---1.16.9, с.~104---105: представимые
 функторы $h_X=\Hom(X,-)$ и теорема $f\mapsto h_f$ --- изоморфизм
 $\Hom_{\mathcal C}(X,Y)\simeq\Hom_{\mathrm{Funct}}(h_X,h_Y)$; функтор
 $h\colon\mathcal C\to\widehat{\mathcal C}$ является погружением
 категории. Именно это и есть категорная форма <<Факта>>:
 $\Hom(A,-)$ определяет объект однозначно.
 Годеман, гл.~I, п.~1.7, с.~25 --- пример того, как морфизмы
 порождают естественные преобразования функторов $\Hom$.
}
\end{bookbox}

%==================================================================
\section[«Прокачивание» через подпредпучок: секции $\Shv{F}^{+}$]
 {<<Прокачивание>> через подпредпучок: секции $\Shv{F}^{+}$}
%==================================================================

Вместо того чтобы каждый раз проверять аксиомы для $\et(\Shv{F})$,
можно <<прокачать>> предпучок внутри заранее готового пучка.

\begin{ass}
Пусть $\Shv{H}$ --- пучок на $X$, а $\Shv{F}\subset\Shv{H}$ ---
подпредпучок: семейства $\Shv{F}(U)\subseteq\Shv{H}(U)$,
согласованные с ограничениями $\rho$.
\end{ass}

\begin{defn}[секции $\Shv{F}^{+}$]
Для открытого $U\subset X$ положим
\begin{equation}\label{eq:Fplus}
 \Shv{F}^{+}(U)=\bigl\{\,s\in\Shv{H}(U)\ \bigm|\ \exists\ \text{покрытие }
 U=\textstyle\bigcup_i U_i\ \text{ и }\ \res^{U_i}_{U}(s)\in\Shv{F}(U_i)
 \ \forall i\bigr\},
\end{equation}
то есть $\Shv{F}^{+}(U)$ --- секции пучка $\Shv{H}$, которые
\emph{локально} являются секциями предпучка $\Shv{F}$.
\end{defn}

\begin{thm}\label{thm:Fplus}
$\Shv{F}^{+}$ --- пучок, и $\Shv{F}^{+}\simeq\Shv{F}^{\#}$.
\end{thm}

Сразу видно, что $\Shv{F}^{+}$ --- пучок: единственность и склейка
наследуются от $\Shv{H}$ (покрытие для двух секций можно объединить,
а для склейки --- взять объединение двух покрытий). Вся работа ---
в изоморфизме.

\subsection{План доказательства}
Вместо прямой проверки используется теорема~\ref{thm:hom} и
Факт~\ref{fact:hom}:
\begin{conditions}
\item построить естественный изоморфизм функторов
 $\Hom(\Shv{F},-)\simeq\Hom(\Shv{F}^{+},-)$ на пучках;
\item изоморфизм $\Hom(\Shv{F},-)\simeq\Hom(\Shv{F}^{\#},-)$ уже есть
 (теорема~\ref{thm:hom});
\item тогда $\Hom(\Shv{F}^{+},-)\simeq\Hom(\Shv{F}^{\#},-)$ естественно
 (композиция изоморфизмов), и <<Факт>> даёт
 $\Shv{F}^{+}\simeq\Shv{F}^{\#}$.
\end{conditions}

\subsection{Построение $\eta_{\Shv{G}}$}
Для пучка $\Shv{G}$ и $\varphi\colon\Shv{F}\to\Shv{G}$ отображение
\[
 \eta_{\Shv{G}}(\varphi)=\bar\varphi\colon\Shv{F}^{+}\to\Shv{G}
\]
строится так.
\begin{conditions}
\item Пусть $s\in\Shv{F}^{+}(U)$. По определению~\eqref{eq:Fplus} есть
 покрытие $U=\bigcup_iU_i$ и $s_i:=\res^{U_i}_{U}(s)\in\Shv{F}(U_i)$.
\item Положим $t_i:=\varphi_{U_i}(s_i)\in\Shv{G}(U_i)$.
\item Семейство $(t_i)$ согласовано: на $U_i\cap U_j$ равенство
 $s_i|=s_j|$ выполняется (оба --- ограничения одного $s$), а $\varphi$
 --- гомоморфизм, поэтому $\varphi(s_i)|=\varphi(s_j)|$.
\item Аксиома (F2) \emph{пучка $\Shv{G}$} даёт единственное
 $t\in\Shv{G}(U)$ с $t|_{U_i}=t_i$. Полагаем $\bar\varphi(s):=t$.
\item Независимость от выбора покрытия: для другого покрытия полученные
 локальные значения на пересечениях совпадают, и снова работает (F2)
 (либо (F1) --- для единственности).
\end{conditions}

Ясно, что $\bar\varphi\circ i=\varphi$ (берём покрытие из одного
множества $U$). Диаграмма~\ref{dgm:eta} фиксирует, что получилось:
$\eta_{\Shv{G}}$ --- естественное по $\Shv{G}$ преобразование
$\Hom(\Shv{F},-)\Rightarrow\Hom(\Shv{F}^{+},-)$.

\begin{diagram}{Слева: $\eta_{\Shv{G}}$ --- естественное по $\Shv{G}$
 преобразование $\Hom(\Shv{F},-)\Rightarrow\Hom(\Shv{F}^{+},-)$;
 справа: $\bar\varphi$ --- единственное продолжение $\varphi$
 до $\Shv{F}^{+}$.\label{dgm:eta}}
\begin{tikzcd}[cd style]
 \Hom(\Shv{F},\Shv{G}_1) \ar[r,"\alpha\circ-"] \ar[d,"\eta_{\Shv{G}_1}"'] &
 \Hom(\Shv{F},\Shv{G}_2) \ar[d,"\eta_{\Shv{G}_2}"] \\
 \Hom(\Shv{F}^{+},\Shv{G}_1) \ar[r,"\alpha\circ-"'] &
 \Hom(\Shv{F}^{+},\Shv{G}_2)
\end{tikzcd}
\qquad
\begin{tikzcd}[cd style,column sep=huge,row sep=huge]
 \Shv{F} \ar[r,"i"] \ar[dr,"\varphi"'] & \Shv{F}^{+} \ar[d,"\bar\varphi"] \\
  & \Shv{G}
\end{tikzcd}
\end{diagram}

Пояснение к диаграмме~\ref{dgm:eta}. Слева --- квадрат
натуральности $\eta_{\Shv{G}}$: домножение морфизма предпучков
на $\alpha\colon\Shv{G}_1\to\Shv{G}_2$ и переход от
$\Hom(\Shv{F},-)$ к $\Hom(\Shv{F}^{+},-)$ коммутируют, т.\,е.\
$\eta$ --- преобразование функторов. Справа --- из чего состоит
компонента $\eta_{\Shv{G}}$: к морфизм $\varphi\colon\Shv{F}\to
\Shv{G}$ единственно продолжается до $\bar\varphi\colon\Shv{F}^{+}
\to\Shv{G}$ с $\bar\varphi\circ i=\varphi$ (треугольник), и
именно $\eta_{\Shv{G}}(\varphi)=\bar\varphi$.

\subsection{Почему $\eta_{\Shv{G}}$ --- изоморфизм}
Остаётся проверить, что каждая компонента $\eta_{\Shv{G}}$ ---
биекция; именно это позволяет применить Факт~\ref{fact:hom}.
\begin{bullets}
\item \emph{Инъективность}: пусть $\eta_{\Shv{G}}(\varphi_1)=
 \eta_{\Shv{G}}(\varphi_2)$. Продолжения совпадают, поэтому после
 композиции с каноническим $i\colon\Shv{F}\to\Shv{F}^{+}$ получаем
 $\varphi_1=\varphi_2$. Иными словами, отображение ограничения
 $\Hom(\Shv{F}^{+},\Shv{G})\to\Hom(\Shv{F},\Shv{G})$ является левым
 обратным к $\eta_{\Shv{G}}$.
\item \emph{Сюръективность}: пусть $\psi\colon\Shv{F}^{+}\to\Shv{G}$.
 Рассмотрим продолжение $\overline{\psi\circ i}$. Для каждого
 $s\in\Shv{F}^{+}(U)$ существует покрытие $U=\bigcup U_i$, на котором
 $s|_{U_i}=i(s_i)$ для $s_i\in\Shv{F}(U_i)$. На каждом $U_i$ оба
 сечения $\overline{\psi\circ i}(s)$ и $\psi(s)$ равны
 $\psi(i(s_i))$. По (F1) для пучка $\Shv{G}$ они равны на $U$.
 Значит, $\eta_{\Shv{G}}(\psi\circ i)=\psi$.
\end{bullets}

Итак $\Hom(\Shv{F},-)\simeq\Hom(\Shv{F}^{+},-)$ естественно, и
Факт~\ref{fact:hom} даёт $\Shv{F}^{+}\simeq\Shv{F}^{\#}$.
\hfill$\square$

\begin{bookbox}
\booksrc{Манин, \S\,2.1.6---2.1.7, с.~112---113: <<секции $\Shv{P}^{+}(U)$
 определяются как некоторые наборы ростков $\prod_{x\in U}\Shv{P}_x$,
 которые в естественном смысле согласованы>>; там же описаны \emph{два}
 процесса <<прокачивания>>: (1) \emph{уменьшить} число сечений,
 отождествляя те, что начинают совпадать на мелком покрытии; (2)
 \emph{увеличить} число сечений, добавляя склеенные из согласованных
 наборов. Теорема-определение 2.1.7: $\Shv{P}^{+}$ --- пучок и
 $(\Shv{P}^{+})_x=\Shv{P}_x$.
\par\smallskip
 Годеман, Замечание 1.2.3, с.~133 --- пучок, порождённый предпучком
 ($\Shv{F}^{\#}=\tilde{\Shv{F}}$); Годеман, п.~3.9, с.~183 ---
 <<локальное равенство» в $\Shv{F}(X)$.
}
\end{bookbox}

%==================================================================
\section{Пример: точные и замкнутые формы}
%==================================================================

\begin{exm}\label{exm:forms}
Пусть $X$ --- гладкое многообразие и $k\geq1$,
$\Shv{F}$ --- предпучок точных ($d$-точных) $k$-форм на $X$,
$\Shv{H}=\Omega^{k}_X$ --- пучок всех $k$-форм. Тогда
$\Shv{F}^{+}$ --- пучок форм, локально точных. По классической лемме
Пуанкаре локально любая замкнутая форма точна, и обратно: точная форма
всегда замкнута ($d^2=0$). Отсюда
\begin{equation}\label{eq:forms}
 \Shv{F}^{+}=\{\text{замкнутые }k\text{-формы}\},\qquad
 \Shv{F}^{\#}\simeq\Shv{F}^{+},
\end{equation}
то есть <<прокачивание>> предпучка точных форм даёт пучок замкнутых
форм: глобально точных форм может быть мало (группы когомологий!),
но локально точных --- столько, сколько хочется.
\end{exm}

\begin{bookbox}
\booksrc{Годеман, гл.~II, \S\,2, с.~156: при доказательстве точности
 используется классическая теорема Пуанкаре (локально любая замкнутая
 форма точна); п.~2.8.1, с.~159 --- примеры пучков $\Omega^p$
 дифференциальных форм и пучков ростков как $\mathcal{O}$-модулей.
 Манин, \S\,2.1.2, с.~110 --- предпучки модулей над предпучком колец.
}
\end{bookbox}

%==================================================================
\section{Ядро, образ и точные последовательности}
%==================================================================

\subsection{Ядро и образ}
\leavevmode\par
\begin{defn}[ядро и образ]
Пусть $f\colon\Shv{F}\to\Shv{G}$ --- гомоморфизм пучков.
\emph{Ядром} $\ker f$ называется пучок, определяемый поэлементно:
\begin{equation}\label{eq:ker}
 (\ker f)(U)=\ker\bigl(\Shv{F}(U)\xrightarrow{f_U}\Shv{G}(U)\bigr),
\end{equation}
ограничения --- обычные. \emph{Образ} определяется по формуле
\begin{equation}\label{eq:im}
 \operatorname{Im}f=\bigl(\operatorname{Im}f\bigr)^{\#},
\end{equation}
где $\operatorname{Im}f$ справа --- предпучок с группами
$\operatorname{Im}\bigl(\Shv{F}(U)\xrightarrow{f_U}\Shv{G}(U)\bigr)$,
а $({-})^{\#}$ --- прокачивание.
\end{defn}

Почему образ \emph{нельзя} оставить поэлементным: множество
$\operatorname{Im}(f_U)$, вообще говоря, не является пучком. Пример:
пусть $\Shv{F}$ --- постоянный предпучок $\mathbb{Z}$, а
$\Shv{G}$ --- пучок, в котором локально ненулевые вещи склеиваются;
образ поэлементно будет <<предпучком, который надо досклеить>>.
Прокачивание~\eqref{eq:im} --- это и есть правильный образ: он
удовлетворяет (F2) по построению.

\begin{rem}[ядро не требует прокачивания]
Ядро --- обратная операция: точность поэлементна сохраняется, потому
что $\Shv{F}\mapsto\Shv{F}(U)$ --- точный функтор (см.~Годемана,
п.~1.9). Поэтому $\ker f$ поэлементно --- уже пучок.
\end{rem}

\subsection{Точные последовательности}
\leavevmode\par
\begin{defn}
Последовательность пучков
$\Shv{F}\xrightarrow{\ f\ }\Shv{G}\xrightarrow{\ g\ }\Shv{H}$
называется \emph{точной}, если $\ker g=\operatorname{Im}f$.
\end{defn}

Диаграмма~\ref{dgm:canon} показывает каноническое разложение
гомоморфизма, из которого и состоит определение.

\begin{diagram}{Каноническое разложение гомоморфизма пучков и
 точность: $\operatorname{Im}f=\ker g$.\label{dgm:canon}}
\begin{tikzcd}[cd style,row sep=large]
0 \ar[r] & \ker f \ar[r,hook] & \Shv{F} \ar[r,"f"'] \ar[d,two heads] &
\Shv{G} \ar[r,"g"] & \Shv{H} \\
 & & \operatorname{Im} f \ar[r,equal] & \ker g \ar[r] \ar[u,hook] & 0
\end{tikzcd}
\end{diagram}

Пояснение к диаграмме~\ref{dgm:canon}. Две вертикали и стрелка
равенства и есть каноническое разложение: $\Shv{F}$ сюръективно
спускается на свой образ $\operatorname{Im}f$ (вниз, <<две
головы>>), образ отождествляется с $\ker g$, а $\ker g$
включается обратно в $\Shv{G}$ (вверх, крючок). Таким образом
$f$ читается как <<сначала спуститься на образ, потом включить>>:
$\Shv{F}\twoheadrightarrow\operatorname{Im}f=\ker g\hookrightarrow
\Shv{G}$. Слева $\ker f\hookrightarrow\Shv{F}$ --- участок, на
котором $f$ равен нулю; полнота (инъективность слева,
сюръективность справа) как раз закодирована нулями по краям.

\subsection{Как проверять точность}
Практическое правило: точность последовательности пучков проверяется
на ростках $\Shv{F}_x\to\Shv{G}_x\to\Shv{H}_x$ для всех $x\in X$.
Эквивалентно, каждое сечение из ядра второго отображения должно
\emph{локально} на покрытии быть образом сечения первого пучка.
Не следует требовать поточечной сюръективности отображений сечений над
каждым открытым множеством: эпиморфизм пучков вообще не обязан быть
сюръективен на глобальных сечениях.

Частый пример: $0\to\Shv{F}\to\Shv{G}$ точна тогда и только тогда,
когда $\Shv{F}\to\Shv{G}$ --- инъекция \emph{в каждом слое}.

\begin{bookbox}
\booksrc{Годеман, гл.~I, п.~1.8, с.~26---29 --- ядро, коядро, образ и
 кообраз в аддитивной категории, определение точной последовательности;
 гл.~II, \S\,2.4---2.5, с.~153---154 --- каноническое разложение
 гомоморфизма пучков и точные последовательности: точность
 последовательности $\mathcal{O}$-модулей эквивалентна точности на
 каждой точке $x$, т.\,е. $\ker(g)(x)=\operatorname{Im}(f)(x)$.
\par\smallskip
 Манин, \S\,2.1.4, с.~110 --- аксиома пучка как точность
 последовательности
 $0\to\Shv{P}(U)\to\prod_i\Shv{P}(U_i)\rightrightarrows
 \prod_{i,j}\Shv{P}(U_i\cap U_j)$;
 \S\,2.9, с.~141 --- обзор предпучков и пучков модулей.
}
\end{bookbox}

%==================================================================
\section{Итог лекции}
%==================================================================

\begin{conditions}
\item Для предпучка $\Shv{F}$ построено локальное накрывающее
 пространство $\pi\colon\et(\Shv{F})\to X$ с
 $\pi^{-1}(x)=\Shv{F}_x=\dirlim_{U\ni x}\Shv{F}(U)$; топология на нём
  --- наиболее тонкая, при которой все сечения непрерывны, и
 $\pi$ --- локальный гомеоморфизм.
\item Пучок $\Shv{F}^{\#}$ определяется как пучок непрерывных сечений
 $\pi$; он удовлетворяет (F1) и (F2): (F1) --- потому что множество
 точек равенства двух сечений открыто, (F2) --- потому что
 согласованные непрерывные выборы ростков склеиваются в
 непрерывный.
\item Для постоянного предпучка $\Shv{F}(U)=A$ имеем
 $\et(\Shv{F})=X\times A$ (дискретное $A$), а $\Shv{F}^{\#}$ ---
 пучок локально постоянных функций.
\item $\Hom(\Shv{F},\Shv{G})\simeq\Hom(\Shv{F}^{\#},\Shv{G})$
 естественно по $\Shv{F}$ и по $\Shv{G}$ при пучке $\Shv{G}$;
 поэтому <<прокачивание>> определено функторно и однозначно
 (правильно: $\Shv{F}\mapsto\Shv{F}^{\#}$ --- левый сопряжённый к
 включению $\Sh\hookrightarrow\PSh$).
\item Естественный изоморфизм $\Hom(\Shv{F},-)\simeq\Hom(\Shv{F}',-)$
 при пучке $\Shv{F}'$ влечёт $\Shv{F}\simeq\Shv{F}'$: следует из
 <<Факта>> о том, что $\Hom(A,-)$ определяет объект категории
 однозначно (лемма Ёнеды).
\item Конструкция $\Shv{F}^{+}$ (секции пучка $\Shv{H}\supset\Shv{F}$,
 локально принадлежащие $\Shv{F}$) даёт то же самое:
 $\Shv{F}^{+}\simeq\Shv{F}^{\#}$; пример --- пучок замкнутых форм из
 предпучка точных.
\item Ядро гомоморфизма пучков определяется поэлементно, образ нужно
 <<прокачивать»: $\operatorname{Im}f=(\operatorname{Im}f)^{\#}$;
  точная последовательность пучков --- это точность на каждом ростке;
  на сечениях над открытым множеством требуется лишь локальное
  представление элемента ядра как образа.
\end{conditions}

\subsection{Что дальше}
Лекция~3, как правило, начинается с того, что все построенные
конструкции (ядро, образ, $\Shv{F}^{+}$) укладываются в \emph{абелеву
категорию} пучков, и переходят к резольвентам и когомологиям. Уже сейчас
видно главное: пучки --- это <<функции>>, которые локально выглядят
просто, а глобальные свойства измеряются коэффициентами при
<<клейке».

%==================================================================
\begin{thebibliography}{9}
\bibitem{god}
 Годеман Р. \emph{Алгебраическая топология и теория пучков}. --- М.:
 Издательство иностранной литературы, 1961. --- Гл.~I, \S\,1.7, 1.8, 1.9;
 гл.~II, \S\,1 (п.~1.1---1.4), \S\,2, \S\,2.4---2.5, п.~3.9.
\bibitem{manin}
 Манин Ю.~И. \emph{Введение в теорию схем и квантовые группы} / под ред.
  Д.~А.~Лейтеса и С.~М.~Львовского. --- М.: МЦНМО, 2012. ---
  \S\,1.10, \S\,1.12; \S\,1.16.5---1.16.9; \S\,2.1.1---2.1.9.
\bibitem{lec}
 Лекция 2 (19.09), рукописный конспект: lecture-02/photo\_1.jpg
 \ldots\ photo\_7.jpg (стр.~9---15 тетради); лекция~1 ---
 lecture-01/.
\end{thebibliography}

\end{document}
